\documentclass[11pt,a4paper]{article}
\usepackage[utf8]{inputenc}
\usepackage{lmodern}
\usepackage{amsmath,amssymb,amsthm,mathtools}
\usepackage{booktabs,array,longtable}
\usepackage{microtype}
\usepackage[left=22mm,right=22mm,top=30mm,bottom=30mm]{geometry}

\usepackage{appendix}

\usepackage{enumitem,needspace}
\usepackage{xcolor,fancyhdr}
\usepackage{xurl}
\usepackage{hyperref}
\usepackage{bookmark}
\definecolor{linkblue}{RGB}{25,62,100}
\hypersetup{colorlinks=true,linkcolor=linkblue,
citecolor=linkblue, urlcolor=linkblue}
\setlength{\emergencystretch}{2em}
\setlength{\parskip}{2pt}
\usepackage{indentfirst}
\setlength{\parindent}{1em}

\numberwithin{equation}{section}
\newtheorem{theorem}{Theorem}[section]
\newtheorem{proposition}[theorem]{Proposition}
\newtheorem{lemma}[theorem]{Lemma}
\newtheorem{conjecture}[theorem]{Conjecture}
\newtheorem{corollary}[theorem]{Corollary}
\theoremstyle{definition}
\newtheorem{definition}[theorem]{Definition}
\theoremstyle{remark}
\newtheorem{remark}[theorem]{Remark}
\newcommand{\Irr}{\operatorname{Irr}}
\newcommand{\IBr}{\operatorname{IBr}}
\newcommand{\Aut}{\operatorname{Aut}}
\newcommand{\Out}{\operatorname{Out}}
\newcommand{\Gal}{\operatorname{Gal}}
\newcommand{\GL}{\operatorname{GL}}
\newcommand{\SL}{\operatorname{SL}}
\DeclareMathOperator{\PSL}{PSL}
\DeclareMathOperator{\PSU}{PSU}
\DeclareMathOperator{\Hom}{Hom}
\DeclareMathOperator{\Inn}{Inn}
\DeclareMathOperator{\St}{St}
\DeclareMathOperator{\bl}{bl}
\newcommand{\Syl}{\operatorname{Syl}}
\newcommand{\Rad}{\operatorname{Rad}}
\newcommand{\dz}{\operatorname{dz}}
\newcommand{\Sym}{\operatorname{Sym}}
\newcommand{\Inf}{\operatorname{Inf}}
\newcommand{\rank}{\operatorname{rank}}
\newcommand{\Nor}{\mathrm{N}}
\newcommand{\Cen}{\mathrm{C}}
\newcommand{\OO}{\mathrm{O}}
\newcommand{\ZZ}{\mathrm{Z}}
\newcommand{\HH}{\mathcal H}
\newcommand{\WW}{\mathcal W}
\newcommand{\FF}{\mathbb F}
\newcommand{\QQ}{\mathbb Q}
\newcommand{\NN}{\mathrm N}
\newcommand{\CC}{\mathrm C}
\newcommand{\et}{\widetilde E}
\newcommand{\AT}{\widetilde A}
\newcommand{\cl}[1]{\bigl[#1\bigr]}

\newcommand{\eps}{\varepsilon}
\newcommand{\defect}{\operatorname{def}}
\newcommand{\wt}[2]{[(#1,#2)]}
\newcommand{\mcite}[2]{\cite[#2]{#1}}

\newcommand{\pair}[1]{\bigl[#1\bigr]}
\newcommand{\Ind}{\operatorname{Ind}}
\newcommand{\varphip}[2]{\varphi^{(#1)}_{#2}}
\setcounter{tocdepth}{2}
\setcounter{secnumdepth}{3}



\title{\bfseries The inductive Navarro-Alperin weight condition\\ for the Suzuki and Ree groups}
\author{Qianhu Zhou}
\date{\scriptsize School of Mathematics and Statistics, Central China Normal University, Wuhan, 430079, China}
\begin{document}
\maketitle
{\noindent\small{\bf Abstract}
In this paper, we give verifications of  the inductive Navarro–Alperin weight condition in \cite{FFZ} for the Suzuki and Ree groups at all primes.}

\medskip\noindent{\small{{{\bf Keywords} Navarro-Alperin weight conjecture; finite groups of Lie type; Galois action}
}}


%\tableofcontents
\section{Introduction}\label{sec:intro}
In the representation theory of finite groups, the so-called local-global conjectures assert a relation between the representation theory of a finite group and that of some of its local subgroups. Fix a prime $p$, and let $G$ be a finite group. We say that a $p$-weight of $G$ is a pair $(R,\varphi)$ if $R$ is a $p$-subgroup of $G$ and $\varphi$ is an irreducible complex character of $\NN_G(R)/R$ such that $\varphi(1)_p=|\NN_G(R)/R|_p$. The Alperin Weight conjecture \cite{Alperin1987} claims that the number of irreducible $p$-Brauer characters of $G$ equals the number of $G$-conjugacy classes of $p$-weights of $G$.
In 2004, Navarro \cite{Navarro2004} refined this conjecture 
and predicted that there is an bijection the set of irreducible $p$-Brauer characters of $G$ and the set of $G$-conjugacy classes of $p$-weights of $G$ such that it can be chosen equivariant under certain Galois automorphisms.
In \cite{FFZ}, this refinement is called as the \textit{Navarro-Alperin weight (NAW) conjecture}, or the \textit{Galois Alperin weight conjecture}(see Conjecture \ref{con2.1} below). In this paper, we are concerned with this conjecture.

Feng, Fu and Zhou \cite{FFZ} adopt the formulation of the NAW conjecture that includes both Galois automorphisms and group automorphisms, and reduce this conjecture to a problem concerning simple groups, yielding a so-called \textit{inductive Navarro-Alperin weight (NAW) condition}(see \cite[Definition 7.1]{FFZ}) that is usually stronger than the original conjecture. The NAW conjecture has been proved in some cases. For example, $p$-solvable groups, finite groups with abelian Sylow $2$-subgroups in \cite{FFZ}. Furthermore, the resulting inductive NAW condition has been verified  for all finite non-abelian simple groups of Lie type in their defining characteristic in \cite{FFZ}, and for the Ree simple groups ${}^2G_2(3^{2m+1})$ with $m\geq 1$ at all primes, and so on.

The main result of this paper is the following theorem.
\begin{theorem}\label{thm:main}
The inductive Navarro-Alperin weight condition is satisfied for  the Suzuki and Ree groups at all primes $p$.
\end{theorem}

The structure of this paper is as follows. 
In Section \ref{section 2}, we introduce some basic notation and a criterion for the inductive NAW condition of some special non-abelian simple groups. In Section \ref{section 3}, we will verify that Suzuki simple groups and Ree simple groups satisfy all hypothesis of Proposition \ref{prop:criterion}. In Section \ref{section 4}, we consider some exceptional cases mentioned at the beginning of Section \ref{section 3}. In Section \ref{section 5}, we will prove Theorem \ref{thm:main}. Finally, in Appendix \ref{app:maps}, we provide some operational steps required to handle an exceptional situation.


\section{Preliminaries}\label{section 2}
We introduce the notation we use throughout this paper. Fix a prime $p$ and a sufficiently large splitting $p$-modular system $(\mathcal K,\mathcal O,F)$ for all finite groups under consideration as in \mcite{FFZ}{Section 1} or \cite[Section 2]{T2013}. Its Galois group over the field \(\QQ_p\) is denoted by \(\HH_p\). On every $p'$-th root of unity, its action is generated by the map \(\sigma_p:\zeta\mapsto\zeta^p\).
The group $\HH_p$ acts on all ordinary and Brauer characters by first evaluating the character and then applying the Galois automorphisms to the character values.

Let $G$ be a finite group. Write $\IBr_p(G)$ and $\Irr(G)$ for the sets of irreducible $p$-Brauer characters and ordinary characters of $G$, respectively. Write $\dz_p(G)$ for the set of $p$-defect zero characters of $G$, that is, for any $\chi\in\dz_p(G)$, we have $\chi(1)_p=|G|_p$.  Let $(R,\varphi)$ be a $p$-weight of $G$. Then $R=\OO_p (\NN_G(R))$ is called the $p$-radical subgroup, and $\varphi \in \dz_p(\NN_G(R)/R)$ is called the weight character. In the following, when mentioning a weight, we often mean this version.  Note that $\varphi^\circ$, restriction of $\varphi$ to all $p'$-elements of $\NN_G(R)/R$, is a projective irreducible Brauer character of $\NN_G(R)/R$. We  can also view $\varphi^\circ$ as an element of $\IBr_p(\NN_S(R))$ via inflation.
Write $\overline{(R,\varphi)}$ for the $G$-conjugacy class containing $(R,\varphi)$ and $\WW_p(G)$ for the set of all such $\overline{(R,\varphi)}$.  

Write $\Aut(G)$ for the automorphism group of $G$. Let $a=(\sigma,\alpha)\in\HH_p\times\Aut(G)$, and let $H$ be 
any subgroup of $G$ and $\chi\in \IBr_p(H)\sqcup\Irr(H)$.
We define $H^a=H^\alpha$ and $\chi^a\in\IBr_p(H^a)\sqcup \Irr(H^a)$
with $\chi^a(h)=\chi(h^{\alpha^{-1}})^\sigma$ for any $h\in H$. 
It is clear that the group $\HH_p\times\Aut(G)$ also acts on $\WW_p(G)$ by separately acting on the two components of the $p$-weight of $G$. With these, we can state the Navarro-Alperin Weight conjecture from \cite{FFZ}.
\begin{conjecture}[{NAW conjecture}]\label{con2.1}
Let $G$ be a finite group $G$ and $\HH_p$  described as above. Then there exists an $(\HH_p\times \Aut(G))$-equivariant bijection from $\IBr_p(G)$ to $\WW_p(G)$.    
\end{conjecture}

We want to show that a simple group $G$ satisfies the inductive NAW condition at the prime $p$ from \cite[Definition 7.1]{FFZ}. Feng, Fu and Zhou in \cite[Proposition 9.1]{FFZ} have given a technical criterion of this inductive condition. Now we only state a its simplified version that implies this inductive  condition for simple groups with trivial Schur multiplier and cyclic outer automorphism group.
\begin{proposition}\label{prop:criterion}
Let \(G\) be a finite non-abelian simple group  with
trivial Schur multiplier and cyclic outer automorphism group. 
Assume that there exists an \((\HH_p\times\Aut(G))\)-equivalent bijection
\begin{equation}\label{equ:bijection}
 \Omega:\IBr_p(G)\longrightarrow\WW_p(G).
\end{equation}
For \(\eta\in\IBr_p(G)\), write
\(\Omega(\eta)=\overline{(R,\varphi)}\).
Assume that there are finite groups $A$, $\widetilde A$ with \(G\unlhd A\unlhd\widetilde A\) such that, under the natural
conjugation maps,
\begin{equation*}
 \widetilde A/\CC_{\widetilde A}(G)\cong\Aut(G)_{\eta^{\HH_p}},\
 A/\CC_{\widetilde A}(G)\cong\Aut(G)_\eta,\ 
 \CC_{\widetilde A}(G)=\ZZ(A).
\end{equation*}
Assume that  \(\eta\) and \(\varphi^\circ\) have extensions $\widehat\eta\in\IBr_p(A)$ and 
$\widehat\varphi\in\IBr_p(\NN_A(R))$
covering the same character of \(\ZZ(A)\), respectively, and that
\begin{equation}\label{extension criterion}
\text{$\widehat \eta$ and $\widehat \varphi$ are $(\HH_p\times\NN_{\widetilde A}(R))_{\varphi}$-invariant}.
\end{equation}
Then \(G\) satisfies the inductive NAW condition at  the prime \(p\).
\end{proposition}


\begin{remark}\label{rem: 2.2}
When  \(G\) has non-trivial Schur multiplier, we use \cite[Proposition 9.1]{FFZ}. Furthermore, the above Proposition can be simplified. For example, when $A$ is a subgroup of $\Aut(G)$ containing $G$, $A/S\cong \NN_A(R)/\NN_{G}(R)$ and  $\CC_{\widetilde A}(G)=\ZZ(A)=1$.  If $A/G$ is a $p'$-group, then as in \cite[Remark 9.2]{FFZ}, the condition \eqref{extension criterion} is satisfied.
\end{remark}
The following proposition is often applied in our subsequent proofs, which is a special case of {\cite[Corollary 10.4]{FFZ}}.
\begin{proposition}\label{cyclic Sylow}
Let $G$ be a finite non-abelian simple group  with trivial Schur multiplier.  Assume that there exists a finite group $A$ such that $G\unlhd A$ and $A$ induces all automorphisms on $G$ by conjugation. 
If $A/G$ is cyclic, $\CC_A(G)=\ZZ(A)$ and $G$ has a cyclic Sylow $p$-subgroup, then the inductive NAW condition holds for $G$ at the prime $p$.
\end{proposition}



\section{Suzuki groups and Ree groups}\label{section 3}
In this Section, we will verify Proposition \ref{prop:criterion} or \cite[Proposition 9.1]{FFZ}  for the very twisted simple groups of Lie type. Several partial results had already been obtained in \cite{FFZ} (see there references in the subsequent proof).

Let $S$ be any of the groups ${}^2B_2(2^{2m+1}), {}^2G_2(3^{2m+1})$, or ${}^2F_4(2^{2m+1})$ for $m\geq 0$ and $\ell$ its the defining characteristic. Denote by $p$ a prime dividing $|S|$ different from $\ell$.  First, the group ${}^2B_2(2)$ is the Frobenius group of order $20$, hence solvable, and we do not have to consider it here. Second, for $m\geq 1$, we have $\Out(S)\cong C_{2m+1}$.
Note that $S$ has trivial Schur multiplier except for ${}^2B_2(8)$, and that $S$ is simple except for \({}^2G_2(3)\) and \({}^2F_4(2)\) that is handled separately  in Section \ref{section 4}.


\subsection{Suzuki and small Ree groups}
\begin{theorem}\label{thm:Suzuki}
Let $S$ be any of the simple groups ${}^2B_2(2^{2m+1})$ or
${}^2G_2(3^{2m+1})$ for $m\geq1$. Then $S$ satisfies the inductive NAW condition at all primes $p$.
\end{theorem}
\begin{proof}
\cite[Proposition 10.9] {FFZ} has verified for  $S={}^2G_2(3^{2m+1})$ at all primes $p$. Thus we suppose that $S={}^2B_2(q^2)$, where $q^2=2^{2m+1}$ with $m\geq 2$. Then the order $|S|=q^4(q^2-1)(q^4+1)$. For $p=2$, then $p$ is the defining characteristic, and we apply \cite[Theorem C]{FFZ}. Let $p$ be an odd prime. Then $p$ divides $q^2-1$ or $q^4+1=(q^2+\sqrt{2}q+1)(q^2-\sqrt{2}q+1)$. Note that for any Sylow $p$-subgroup of $P$, 
there exist a Sylow $\Phi^{(p)}$-torus $\bf T$ of ${\bf S}$ such that $P\leq T={\bf T}^F$.  Also, by \cite[Section 16]{I2007}, these relevant tori $T$ are cyclic. So every Sylow $p$-subgroup of $S$ is cyclic, and thus we apply Proposition \ref{cyclic Sylow} to $A=\Aut(S)$.

For the group ${}^2B_2(8)$, its the Schur multiplier is $C_2\times C_2$. So we verify \cite[Proposition 9.1]{FFZ}. Let $p$ be odd. Then the universal $p'$-covering group of ${}^2B_2(8)$ is  $L=2^2.{}^2B_2(8)$.  Note that $\Out(L)\cong L\rtimes C_3$. Also, note that every Sylow $p$-subgroup of $S$ is cyclic. Then we apply \cite[Corollary 10.4]{FFZ}. We complete the proof.
\end{proof}
\subsection{Large Ree groups}
Let $S={}^2F_4(q^2)$ with $q^2=2^{2m+1}$ and $m\geq 1$.  Then the order formula is
\[
|S|=q^{24}(q^2-1)^2(q^2+1)^2(q^4+1)^2(q^4-q^2+1)(q^8-q^4+1),
\]
and $\Aut(S)\cong S\rtimes C_{q^2}$. Moreover, the center and Schur multiplier of $S$ are trivial, and write $D=\Out(S)\cong C_{q^2}$, which consists only of  field automorphisms. 

We will verify all conditions of Proposition \ref{prop:criterion} for $S$ at every prime $p$. First, we consider the case where a Sylow $p$-subgroup is cyclic.  
\begin{lemma}\label{lem:3.2}
Assume that $p$ is odd and $p\nmid q^8-1$. Then $S$ satisfies the inductive NAW condition at the prime $p$.
\end{lemma}
\begin{proof}
Since $p$ is odd and $p\nmid q^8-1=(q^4+1)(q^2+1)(q^2-1)$, we know that $p\geq5$ and $p$ divides $q^4-q^2+1$ or $q^8-q^4+1$. In this case, any sylow $p$-subgroup of $S$ is contained in a sylow $\Phi^{(p)}$-torus. By \cite[TABLE I]{An}, there exists such cyclic Sylow $\Phi^{(p)}$-torus of order $q^4-q^2+1$. Also, $q^8-q^4+1=u_+u_-$, where $u_\pm=q^4+q^2+1\pm \sqrt{2}q(q^2+1)$. Similarly, by \cite[TABLE I]{An}, both $u_\pm$ occur as orders of cyclic Sylow tori.  Their greatest common divisor is $1$ because a common odd divisor would divide both $q^4+q^2+1$ and $\sqrt{2}q(q^2+1)$, which are coprime. Thus there exists a cyclic torus contains a Sylow $p$-subgroup in every case under consideration. Apply Proposition \ref{cyclic Sylow} to $A=\Aut(S)$, we complete the proof. 
\end{proof}
Therefore, it suffices to consider the following two cases in subsections \ref{p>=5} and \ref{p=3} 
\begin{equation}\label{eq:remaining}
 p=3,\ \text{or}\ p\geq5,\ p\mid(q^2-1)(q^2+1)(q^2-\sqrt{2}q+1)(q^2+\sqrt{2}q+1).
\end{equation}

\subsubsection{An equivariant bijection for \texorpdfstring{$p\geq5$}{p >= 5}}\label{p>=5}
In this Subsection, we assume that  $p\geq5$ and $p\mid(q^2-1)(q^2+1)(q^2-\sqrt{2}q+1)(q^2+\sqrt{2}q+1)$.
Then it follows from \cite{Johansson} that any Sylow $p$-subgroup  of $S$ is abelian of rank-two.  Also, by \cite[Corollary~4.3 and the decomposition tables in Appendix A]{Himstedt} and together with \cite{GH}, the ordinary basic set $\mathcal{B}:=\bigcup_{s}\mathcal{E}(S,s)$,
is unitriangular, where $s$ runs over the set of semisimple $p'$-elements. Moreover, it is stable under the actions of
required field automorphisms and Galois automorphisms. Therefore, it follows from \cite[Lemma~10.7]{FFZ} that there exists  an $(\HH_p\times D)$-equivalent bijection 
\begin{equation}\label{eq:basic}
\mathcal{B}\longleftrightarrow\IBr_p(S).
\end{equation}
Note that this bijection \eqref{eq:basic} is independent of the value of $p$, so it can also be used when $p=3$.

We will construct an bijection between $\mathcal{B}$ and $\WW_p(S)$, and prove that it's $(\HH_p\times D)$-equivalent.
Next, based on the different of $p$-radical subgroups, we will discuss them separately.



\begin{lemma}[{\cite[(2B)]{An}}]\label{p3:prop:radical}
There are exactly two non-trivial \(S\)-conjugacy classes of \(p\)-radical subgroups, represented by $R$ and $P$, where $R$ is cyclic and $P$ is non-cyclic abelian.
\end{lemma}
\subparagraph{(i) Non-cyclic abelian \(p\)-radical subgroups.}
Let $P$ be non-cyclic abelian \(p\)-radical subgroup of $S$. Then by \cite[(2B)-(b)]{An}, $P$ is a Sylow $p$-subgroup of $S$. Moreover, $\CC_S(P)=T={\bf T}^F$ and $P=\OO_p(T)$, where ${\bf T}$ is a Sylow $\Phi^{(p)}$-torus of $\bf S$. The normalizer of the torus ${\bf T}$ as following Table \ref{tab:1} (\cite[TABLE I]{An}). 
Here, $D_{16}$ is a dihedral group of order $16$, and $G_8$ denotes the Shephard--Todd complex reflection group of order $96$.
\begin{table}[ht]
\centering
\caption{Sylow tori and normalizers}\label{tab:1}
\begin{tabular}{@{}lllrr@{}}
\toprule
Prime condition &Torus $T$ & $W=\NN_S({\bf T})/T$\\
\midrule
$p\mid q^2-1$ & $C_{q^2-1}\times C_{q^2-1}$ & $D_{16}$\\
$p\mid q^2+1$ & $C_{q^2+1}\times C_{q^2+1}$ & $\GL_2(3)$\\
$p\mid q^2\pm\sqrt{2}q+1$ & $C_{q^2\pm\sqrt{2}q+1}\times C_{q^2\pm\sqrt{2}q+1}$ & $G_8$\\
\bottomrule
\end{tabular}
\end{table}
It is not difficult to verify that $\NN_S(P)=\NN_S({\bf T})\cong T\rtimes W$, where $W$ is a $p'$-group. In this case, we have the following result. We write $\Irr_{p'}(S)$ for the set of all $p'$-characters(their degrees are not divisible by $p'$), and $\Aut(S)_P$ for the subset of $\Aut(G)$ stabilizing $P$.
\begin{proposition}\label{bij: rank-two}
Let $P$ be non-cyclic abelian \(p\)-radical subgroup of $S$. Then there exists an $(\HH_p\times \Aut(S)_P)$-equivalent bijection
\begin{equation}
\Omega_P:\Irr_{p'}(S)\cap \mathcal{B}\longrightarrow\dz_p(\NN_S(P)/P) 
\end{equation}
Moreover, for every \(\chi\in\Irr_{p'}(S)\cap \mathcal{B}\), if we replace $\Omega_P(\chi)$ by  $\overline{(P,\Omega_P(\chi))}$, then the resulting bijection is $(\HH_p\times \Aut(S))$-equivalent.
\end{proposition}
\begin{proof}
Although \(S^*\cong S\), we retain the star to distinguish dual semisimple elements from local character. Fix the standard parametrization
$T^*\longrightarrow\Irr(T), s\mapsto\theta_s$.
For any \(s\in T^*\), let $N=\NN_S(P)$ $N_s=\NN_S(P)_{\theta_s}$ and $W_s=\NN_S({\bf T})_{\theta_s}/T$. Through the dual identification, \(W_s\) is identified with the relative Weyl group associated with the chosen torus inside \(\CC_{S^*}(s)\).

Note that $N=\NN_S(P)=\NN_S({\bf T})$, by \cite[Lemma 3.3]{Johansson}, we have an $(\HH_p\times \Aut(S)_P)$-equivalent  extension map 
\[
\Lambda:\Irr(T)\longrightarrow\bigcup_{T\leq I\leq N}\Irr(I)
\]
with $\Lambda(\theta_s)\in\Irr(N_s)$ and $\Lambda(\theta_s)|_T=\theta_s$.
For any $\lambda\in\Irr(W_s)$, define 
\[
\psi_{s,\lambda}=\Ind_{N_s}^{N}(\Lambda(\theta_s)\Inf_{W_s}^{N_s}(\lambda)).
\]
Clifford theory gives all irreducible characters of \(N\) in this way. Since \(p\nmid |W|\), the degrees
$\psi_{s,\lambda}(1)=|N:N_s|\lambda(1)$ are prime to \(p\). Consequently, we have $\Irr(N)=\Irr_{p'}(N)$. 
For any $\sigma\alpha\in\HH_p\times \Aut(S)_P$, by \cite[Proposition 3.4]{Johansson} and its proof,   $(\psi_{s,\lambda})^{\sigma\alpha}=
\psi_{s^{\sigma\alpha},\,\lambda^{\sigma\alpha}}$.

We claim that $P\leq\ker(\psi_{s,\lambda})$ if and only if $s$ is $p'$-element.  Note that 
$\psi_{s,\lambda}|_T=\lambda(1)
\sum_{\theta_i}\theta_i$, where $\theta_i$ runs over the set of \(N\)-conjugates of \(\theta_s\). Since \(P\unlhd N\), all these conjugates are trivial on \(P\)
if \(\theta_s\) is trivial on \(P\). Conversely, if \(\psi_{s,\lambda}|_P\) is a multiple of \(1_P\), then every constituent in the restriction must be
trivial on \(P\). Therefore, $P\leq\ker(\psi_{s,\lambda})$ if and only if 
$P\leq\ker(\theta_s)$. Since $P=\OO_p(T)$,
any  \(\theta_s\in\Irr(T)\) is trivial on \(P\) if and only if its order is prime to \(p\), as claim.

Let $J_s:\mathcal{E}(S,s)\longrightarrow\mathcal{E}(\CC_{S^*}(s),1)$ be the Jordan decomposition and 
$\mathcal I_s:\Irr(W_s)\longrightarrow
\mathcal{E}(\CC_{S^*}(s),{\bf T}^*,1)$ the relative Weyl-group correspondence(see \cite[Sections 3.1-3.2]{Johansson}).  Define $\chi_{s,\lambda}=J_s^{-1}(\mathcal I_s(\lambda))$. Using \cite[Theorem B and Proposition~3.2]{Johansson}, for any $\sigma\alpha\in\HH_p\times \Aut(S)_P$,  we have  $(\chi_{s,\lambda})^{\sigma\alpha}=
\chi_{s^{\sigma\alpha},\,\lambda^{\sigma\alpha}}$.

Consider all the parametrization classes above  \(\Irr_{p'}(S)\). Restricting to semisimple elements of order prime to \(p\) gives exactly the subset $\Irr_{p'}(S)\cap \mathcal{B}$. The same parametrization give exactly \(\Irr_{p'}(N/P)\) on the local side. 
Since \(P\) is a Sylow \(p\)-subgroup of \(N\), 
\(N/P\) is a \(p'\)-group. Thus $\dz_p(N/P)=\Irr_{p'}(N/P)$. Let \(\bar\psi_{s,\lambda}\) denote the character of \(N/P\) whose inflation is \(\psi_{s,\lambda}\), and define $\Omega_P(\chi_{s,\lambda})=\overline\psi_{s,\lambda}$. Then by the above process, $\Omega_P$ is
$(\HH_p\times \Aut(S)_P)$-equivalent.

Finally, define 
$\widetilde\Omega_P(\chi)=\overline{(P,\Omega_P(\chi))}$. Let \(\alpha\in\Aut(S)\). There is an inner automorphism \(c_g\) such that $\beta=c_g\circ\alpha\in\Aut(S)_P$. Note that \(\alpha\) and \(\beta\) induce the same actions on these two sets. Thus $\widetilde \Omega_P$ is $(\HH_p\times \Aut(S))$-equivalent, we complete the proof.
\end{proof}

\subparagraph{(ii) Cyclic $p$-radical subgroups.} 
Let $\mathbf{M}$ be a $\Phi$-split Levi subgroup of $\bf S$. Write $M={\bf M}^F$ and $Z=\ZZ(M)$. Let $R$ be a non-trivial cyclic $p$-radical subgroup of $S$. Then as in \cite[Section 3]{Malle}, $R=Z_p$, $M=\CC_{S}(R)$ and $\NN_{S}(R)=\NN_{S}({\bf M})$. 
In fact,  $R=\OO_p(\CC_S(R))$ by \cite[(2B)-(a)]{An},
and $M=Z\times K$. Moreover, by \cite[TABLE I]{An}, if $p|q^2\pm 1$, then $\NN_S(R)/\CC_S(R)\cong C_2$. If  $p|q^2\pm\sqrt{2}q+1$, then $\NN_S(R)/\CC_S(R)\cong C_4$.
The centralizer and normalizer of $R$ as in following Table 2. 
\begin{table}[ht]
\centering
\caption{Centralizers and normalizers of cyclic $p$-radical subgroups}\label{tab:2}
\begin{tabular}{@{}lllrr@{}}
\toprule
Prime condition & $\CC_S(R)=Z\times K$ & $\NN_S(R)$\\
\midrule
$p\mid q^2-1$ & $C_{q^2-1}\times\SL_2(q^2)$ & $(C_{q^2-1}\rtimes C_2)\times\SL_2(q^2)$\\
$p\mid q^2-1$ & $C_{q^2-1}\times{}^2B_2(q^2)$ & $(C_{q^2-1}\rtimes C_2)\times{}^2B_2(q^2)$\\
$p\mid q^2+1$ & $C_{q^2+1}\times\SL_2(q^2)$ & $(C_{q^2+1}\rtimes C_2)\times\SL_2(q^2)$\\
$p\mid q^2-\sqrt{2}q+1$ & $C_{q^2-\sqrt{2}q+1}\times{}^2B_2(q^2)$ & $(C_{q^2-\sqrt{2}q+1}\rtimes C_4)\times{}^2B_2(q^2)$\\
$p\mid q^2+\sqrt{2}q+1$ & $C_{q^2+\sqrt{2}q+1}\times{}^2B_2(q^2)$ & $(C_{q^2+\sqrt{2}q+1}\rtimes C_4)\times{}^2B_2(q^2)$\\
\bottomrule
\end{tabular}
\end{table}
Here, these structures are determined by the classification of centralizers of semisimple elements and the structure of  normalizers of tori (see \cite[TABLE I]{An} or \cite[Tables~B.9--B.10]{HH}). For example, in \cite[TABLE I]{An},  from the first row to the end, 
%we have $\CC_S(R)=\CC_S(t_i)$($i=2,1,5,7,9$), respectively.
corresponding to the centralizer $\CC_{S}(t_i)$($i=2,1,5,7,9$) contains the factor $A_1$, ${}^2B_2$, $A_1$, ${}^2B_2$, ${}^2B_2$, respectively. 

For completeness, we need to show that there is no proper radical subgroup whose centralizer is a maximal torus. To see this, suppose that $\CC_S(R)=T$ is a maximal torus. Then $T_p$ is normal in $\NN_S(R)$, so the radicality forces $R=T_p$, which is a Sylow subgroup, and this is contradictory. 
Note that the exceptional semisimple centralizer $\CC_S(t_4)$ of type ${}^2A_2$ in \cite[TABLE I]{An} occurs at an element of order $3$, and does not add a $p$-radical subgroup for $p\geq5$.

Write $E=\NN_S(R)/\CC_S(R)$. Then $E\cong C_2$ or $C_4$. Since \(Z\) is cyclic, $Z=R\times Z_{p'}$,
and then
\[
\NN_S(R)/R=(Z_{p'}\rtimes E)\times K.
\]
As in \cite[Sections 3 and 5]{Malle}, by Clifford theory, any element of $\dz_p(\NN_S(R)/R)$ is parametrized by a triple \((\theta,\kappa,\epsilon)\), where \(\theta\in\Irr(Z_{p'})\), \(\kappa\in\dz_p(K)\) and \(\epsilon\in\Irr(E_\theta)\). In deed, the character \(\theta\) can be extended to $E_\theta$ by defining
$\widehat\theta(ze)=\theta(z)$ for any $z\in Z_{p'}$ and $e\in E_\theta$. Let
\begin{equation}\label{eq:rankone}
 \psi_{\theta,\kappa,\epsilon}=
 \Ind_{(Z_{p'}\rtimes E_\theta)\times K}^{(Z_{p'}\rtimes E)\times K}
((\widehat\theta\epsilon)\otimes\kappa).
\end{equation}
Note that \(E_\theta\) is cyclic and \(\epsilon(1)=1\), we have
$\psi_{\theta,\kappa,\epsilon}(1)
=|E:E_\theta|\kappa(1)$. Since \(|E:E_\theta|\) is a \(p'\)-number, it follows that
$\psi_{\theta,\kappa,\epsilon}(1)_p=\kappa(1)_p$.
On the other hand, $|\NN_S(R)/R|_p=|K|_p$.
Then $\psi_{\theta,\kappa,\epsilon}\in\dz_p(\NN_S(R)/R)$ if and only if  $\kappa\in\operatorname{dz}_p(K)$.

We claim that $E_\theta=E$ for $\theta=1$ and $E_\theta=1$ otherwise. Suppose that  \(E=C_2\).  Then the action of \(C_2\) on \(Z_{p'}\) is inversion. Since \(|Z_{p'}|\) is odd, $E_\theta=C_2$ for $\theta=1$ and $E_\theta=1$ otherwise. Suppose that 
\(E=C_4\), the square of the action  is inversion. Accordingly, we may choose a generator so that the action is given by the \(q^2\)-th power.
Note that $(q^2+\sqrt{2}q+1)(q^2-\sqrt{2}q+1)=q^4+1$  and $q^4\equiv-1\pmod{q^2\pm\sqrt{2}q+1}$.
If a non-trivial subgroup stabilizes \(\theta\), then it contains the unique involution of \(C_4\), and so we again have \(\theta=\theta^{-1}\). Thus
$E_\theta=C_4$ for $\theta=1$ and $E_\theta=1$ otherwise, as claimed.

We can choose an odd-order cyclic subgroup $D_0\leq D$ representing the field automorphism such that it preserves $Z_{p'}$, $K$ and satisfies $[D_0,E]=1$ as follows. Write $I=Z_{p'} \rtimes E$ and take a lift $D_1$ of the field-action preserving $I$. Since $E$ is the Sylow $2$-subgroup of $I$, the Frattini argument yields $I \rtimes D_1=I\NN_{I\rtimes D_1}(E)$.
Note that $\NN_I(E)=\CC_{Z_{p'}}(E)\rtimes E$.
Since $E=C_2$ or $C_4$, and $Z_{p'}$ has odd order, $\CC_{Z_{p'}}(E)=1$, and then $\NN_I(E)=E$. Thus
$$1 \longrightarrow E \longrightarrow \NN_{I\rtimes D_1}(E) \longrightarrow D_1 \longrightarrow 1.$$
By the Schur-Zassenhaus theorem, we may select an odd-order complement $D_0$ isomorphic to $D_1$ inside this normalizer. Note that $\Aut(E)=1$ or $C_2$. So the conjugation action of $D_0$ on $E$ must be trivial. Therefore, $D_0$ acts on all sets of their characters.
\begin{lemma}\label{lem:3.5}
Let $\sigma\in\HH_p$ and $\alpha\in D_0$ normalize $R$. Then $\psi_{\theta,\kappa,\epsilon}^{\sigma\alpha}=\psi_{\theta^{\sigma\alpha},\kappa^{\sigma\alpha},\epsilon^{\sigma}}$.
\end{lemma}
\begin{proof}
Since $|Z_{p'}|$ is odd and $E_\theta$ has order $1$ or $4$, we know that $\widehat \theta$ is the unique linear extension of odd degree. Since $E_\theta=E$ or $E_\theta=1$, $D_0$ acts trivially on $E_\theta$, and then any $\epsilon$ is $D_0$-fixed. For any $\sigma\in\HH_p$ and $\alpha\in D_0$, $\widehat \theta^{\sigma\alpha}$ is an extension of $\theta^{\sigma\alpha}$. By uniqueness of the extension, we have $\widehat \theta^{\sigma\alpha}=\widehat{ \theta^{\sigma\alpha}}$. Both induction and tensor products of characters commute with the action of $\sigma\alpha$. We complete the proof. 
\end{proof}





\begin{lemma}\label{defect zero characters of $K$}
Let $K$ be as in Table \ref{tab:2}. Then its unipotent characters as following Table. Here, $\St$ denotes the Steinberg character.
\begin{table}[ht]
\centering
\begin{tabular}{lll}
\toprule
\(K\) & The unipotent \(\kappa\) & degrees \\
\midrule
\(\SL_2(q^2)\) &  $1,\St$ & $1, q^4$ \\
\({}^2B_2(q^2)\)& $1,\St, W^+,W^-$ & $1,q^4, (\sqrt{q^2/2})(q^2-1), (\sqrt{q^2/2})(q^2-1)$ \\
\bottomrule
\end{tabular}
\end{table}
Moreover, for any \(\sigma\alpha\in\HH_p\times D_0\),  
$\St^{\sigma\alpha}=\St$, $(W^\pm)^\alpha=W^\pm$ and 
\[(W^\pm)^\sigma=
		\begin{cases}
			W^\pm,&i^\sigma=i,\\
			W^\mp,&i^\sigma=-i,
		\end{cases}
        \]
where $i$ denotes a primitive fourth root of unity.
\end{lemma}
\begin{proof}
Suppose that \(K=\SL_2(q^2)\). Then its two unipotent characters are $1$ and $\St$ with degrees $1$ and $q^4$(see \cite[Section 4]{W13}).
Suppose that $K={}^2B_2(q^2)$. Then there are four unipotent characters $1,\St,W^+,W^-$ with degrees $1,q^4,(\sqrt{q^2/2})(q^2-1),(\sqrt{q^2/2})(q^2-1)$ (see \cite[Section 2]{Ghaffarzadeh}).  

Clearly, $\St$ is fixed because Galois automorphisms and field automorphisms preserve degrees of characters. For \(W^+\) and \(W^-\), they are individually fixed by field automorphisms by \cite[Lemma 2.2]{Ghaffarzadeh}. By \cite[Proof of Theorem B]{Johansson}, their Galois actions are determined by the action on \(i\). They are fixed if \(i^\sigma=i\), and interchanged if \(i^\sigma=-i\). 
\end{proof}

Let $B$ be a $p$-block of $S$ parametrized by the cuspidal pair $({\bf M},\lambda)$ of $\bf S$(see \cite[Section 3]{Malle}). Let $R$ be \(p\)-radical subgroup of $S$. Then by \cite[Corollary 3.2 and Remark 2.5]{Malle}, we have $R=Z_p=\OO_p(M)$ and then $M/R=Z_{p'}\times K$.
\begin{proposition}\label{bij: rank-one}
Let $R$ be cyclic \(p\)-radical subgroup of $S$. Then there exist an $(\HH_p\times \Aut(S))_{R,\lambda}$-equivalent bijection
\begin{equation}
\Omega_R:
\mathcal{E}(S,s)\cap\Irr(B)\longrightarrow \dz_p(\NN_S(R)/R|\lambda) 
\end{equation}
Moreover, for every \(\chi\in\mathcal{E}(S,s)\cap\Irr(B)\), if we replace $\Omega_R(\chi)$ by $\overline{(R,\Omega_R(\chi))}$, then the resulting bijection is $(\HH_p\times \Aut(S))$-equivalent.
\end{proposition}
\begin{proof}
For the cuspidal pair $({\bf M},\lambda)$. Recall that $W_S({\bf M},\lambda)=\NN_{S}({\bf M},\lambda)/M$, $\NN_{S}({\bf M})=\NN_S(R)$ and $M=\CC_S(R)$. So $W_S({\bf M},\lambda)\cong E_\lambda$. Thus the parameterization chain in \cite[Propositions~3.5--3.6 and Corollary~3.7]{Malle} is
\begin{equation}\label{eq:chain}
\mathcal{E}(S,s)\cap\Irr(B)\longleftrightarrow\Irr(E_\lambda)
 \longleftrightarrow\dz_p(\NN_S(R)/R\mid\lambda),\ 
\chi_{\lambda, \epsilon}\leftrightarrow \epsilon\leftrightarrow \psi_{\lambda,\epsilon}.
\end{equation}
Here, $\lambda=\theta\otimes\kappa\in\dz_p(M/R)$; $\psi_{\lambda,\epsilon}$ is identified with $\psi_{\theta,\kappa,\epsilon}$ in \eqref{eq:rankone} since $E_\lambda=E_\theta$; The character $\chi_{\lambda,\epsilon}$ denotes an irreducible constituent of the generalized Harish–Chandra induction $R_M^G(\lambda)$.



We do not assume that an arbitrary generalized Harish--Chandra bijection is automatically Galois-equivariant. Instead choose a representative of each $(\HH_p\times D_0)$-orbit of cuspidal pair $({\bf M},\lambda)$, we choose the correspondence there, and transport it. The following stabilizer check makes the construction well-defined. 

Since $E=C_2$ or $C_4$, we have $E_\lambda=1$, $C_2$ or $C_4$. Suppose that $E_\lambda=1$. Then both sets in \eqref{eq:chain} are singletons. This is trivial.  
Suppose that $E_\lambda\cong C_2$. Then
we assume that $\Irr(E_\lambda)=\{1,sgn\}$, where 
$sgn$ is the symbol character. 
Then the two global characters $\chi_{\lambda,1}, \chi_{\lambda,sgn}\in\mathcal{E}(S,s)$. By Jordan decomposition and Table 2, they correspond to the unipotent characters $1$ and $\St$ of $\CC_S(s)$, respectively, where $\CC_S(s)$ is of type $A_1$ or ${}^2B_2$.
Since $1$ and $\St$ have  different degrees, $\chi_{\lambda,1}$ and $\chi_{\lambda,sgn}$ have different degrees. Galois automorphisms and field automorphisms preserve  degrees of characters, so both characters are fixed.   On the other hand, the local characters are $\psi_{\lambda, 1}$ and $\psi_{\lambda,sgn}$. Since  $1$ and $sgn$ are rational, and together with Lemma \ref{lem:3.5}, $\psi_{\lambda, 1}$ and $\psi_{\lambda,sgn}$ are also fixed. 

Suppose that $E_\lambda=C_4$. Let $c$ generate $C_4$.
Then all characters of $C_4$ are $\eta_j(c)=i^j$,
where $i$ is a primitive fourth root of unity and 
$j=0,1,2,3$. Thus we write
$\Irr(E_\lambda)=\{\eta_j|j=0,1,2,3\}$.
Note that \(p\mid q^4+1\) and \(p\) is odd, so by \cite[Lemma 3.1(d)]{Johansson}, $p\equiv 1\pmod 4$.  For any \(\sigma\in\HH_p\), we have $\sigma(i)=i^{p^r}=i$. Thus every $\eta_j$ is \(\HH_p\)-fixed. Clearly, field automorphisms in $D_0$ also fix them. By Lemma \ref{lem:3.5}, these characters $\psi_{\lambda,\eta_j}$ are fixed. 
	
It remains to prove the global part. In the case,
the semisimple centralizer $\CC_S(t_i)$  is of type 
${}^2B_2(q^2)$ for $i=7,9$(in \cite[TABLE I]{An}).
By Lemma \ref{defect zero characters of $K$}, $\CC_S(t_i)$ has only four unipotent characters $1,\St,W^+,W^-$. Clearly, \(1\) and \(\St\) is fixed. Also, since $\sigma(i)=i$, by Lemma \ref{defect zero characters of $K$} again,  $W^{\pm}$ is also fixed. Therefore, all characters $\chi_{1,\eta_j}(j=0,1,2,3)$ are fixed.

Form the above process, for any $\chi_{\lambda,\epsilon}\in\mathcal{E}(S,s)\cap \Irr(B)$, 
we define $\Omega_R(\chi_{\lambda,\epsilon})=\psi_{\lambda,\epsilon}$. Since $\Aut(S)=S\rtimes D$ and $M=\CC_S(R)$,  $\Omega_R$ can be chosen as  $(\HH_p\times\Aut(S))_{R,\lambda}$-equivalent.
Finally, define 
$\widetilde\Omega_R(\chi)=\overline{(R,\Omega_R(\chi))}$. Let \(\alpha\in\Aut(S)\). There is an inner automorphism \(c_g\) such that $\beta=c_g\circ\alpha\in\Aut(S)_{R,\lambda}$. Note that \(\alpha\) and \(\beta\) induce the same actions on these two sets. Thus $\widetilde \Omega_R$ is $(\HH_p\times \Aut(S))$-equivalent, we complete the proof.
\end{proof}
For any $\chi\in\dz_p(S)$, the assignment $\chi^\circ\mapsto\overline{(1,\chi)}$ is obviously  equivariant. Together with Proposition \ref{bij: rank-two}, Proposition \ref{bij: rank-one} and the bijection \eqref{eq:basic}, we obtain
\begin{theorem}\label{equ:bijection when p>=5}
Keep the notation as above. Then there exists an $(\HH_p\times\Aut(S))$-equivalent bijection
$\Omega:\IBr_p(S)\longrightarrow\WW_p(S)$.
\end{theorem}

\subsubsection{Compatible extensions for \texorpdfstring{$p\geq5$}{p >= 5}}
We keep the groups $S$, $D$ and $\HH_p$ as above. We will verify that the bijection in Theorem \ref{equ:bijection when p>=5} satisfies the condition \eqref{extension criterion}. For global characters, since $\mathcal{B}$ is $(\HH_p\times \Aut(S))$-invariant, by \cite[Lemma 10.4(2)]{FFZ}, we only need to consider the ordinary characters. Note that \(\Irr(S)\) contains a large number of real-valued characters, and only a portion of irreducible characters are not real-valued(see \cite{Himstedt}).

The real-valued characters and unipotent characters are handled by the following two Lemmas, respectively.
\begin{lemma}\label{lem:real}
Let $S\unlhd G$ with $G/S$ cyclic with odd order. If $\theta\in\Irr(S)$ is real-valued, then it has a unique real-valued extension to $G_\theta$. Moreover, this extension is invariant under the actions of group and Galois automorphisms stabilizing $\theta$ and preserving $G$ and $S$. 
\end{lemma}
\begin{proof}
The cyclic quotient gives the extendibility. Choose an extension $\chi$ and let $X=\Irr(G_\theta/S)$. Then $X$ has odd order, and acts regularly on the extensions of $\theta$. Since $\bar\theta=\theta$, write $\bar\chi=\mu\chi$ with $\mu\in X$. Then $\overline{\lambda\chi}=\lambda^{-1}\mu\chi$.
So $\overline{\lambda\chi}=\lambda\chi$ if and only if $\lambda^2=\mu$.
Thus the squaring map on $X$ is bijective. This proves existence and uniqueness. Galois automorphisms and group automorphisms commute with $-$, so uniqueness proves the invariance of actions, as required.
\end{proof}

\begin{lemma}[{\cite[Proposition~4.3]{Johansson}}]\label{lem:unipotent}
If $\theta\in\Irr(S)$ is unipotent, then it exists an $(\HH_p)_\theta$-invariant extension to $S\rtimes D$.  
\end{lemma}

For the remaining characters of $S$, with the CHEVIE numbering used in \cite[Table~A.2]{Himstedt} and \cite[Appendix~B]{HH}, the families of non-unipotent and non-real valued characters that need to be processed separately are
\begin{equation}\label{eq:six}
\chi_{23},\chi_{24},\chi_{43},\chi_{44},\chi_{47},\chi_{48}.
\end{equation}
%with numbers $\frac{1}{2}(q^2-2), \frac{1}{2}(q^2-2), \frac{1}{4}(q^2-\sqrt{2}q),\frac{1}{4}(q^2-\sqrt{2}q),\frac{1}{4}(q^2+\sqrt{2}q),\frac{1}{4}(q^2+\sqrt{2}q)$.
Note that in \cite[Table~B.12]{HH}, the values of the last four contain non-zero imaginary parts. 

\begin{lemma}\label{lem: The split Suzuki Levi families}
If $\chi$ is any of characters in the families 
$\chi_{23},\chi_{24}$, then it has an $(\HH_p\times D)_\chi$-invariant extension to $S\rtimes D_\chi$.  
\end{lemma}
\begin{proof}
Let ${\bf J=MU}$ be the standard parabolic subgroup with the split Levi subgroup ${\bf M}^F=M\cong C_{q^2-1}\times{}^2B_2(q^2)$. 
For a non-trivial $\theta\in\Irr(C_{q^2-1})$ and a cuspidal unipotent character $W\in\{W^+,W^-\}$ as in Lemma \ref{defect zero characters of $K$}, we consider the Harish--Chandra induction
\begin{equation}\label{eq:hc}
 \chi=\Ind_{J}^S\Inf_M^{J}(\theta\otimes W).
\end{equation}
By \cite[Proposition~10.4(1), its proof, and Table~1]{HHM},
$\chi$ is irreducible, and the induction gives the families $\chi_{23}$ or$\chi_{24}$.

Let $d\in D$. Then we claim that there exists a $d$-fixed involution $n\in S$ such that 
\begin{equation}\label{eq: act}
(\theta\otimes W)^n=\theta^{-1}\otimes W.   
\end{equation}
To see this,
we prove that there exists the element $n$ satisfies $z^n=z^{-1}$ for any $z\in C_{q^2-1}$ and $[n,{}^2B_2(q^2)]=1$.
Let $w=w_0w_1$, where $w_0$ is the length element in root system of \({}^2F_4\) and the length element \(w_1\) in root subsystem $\Phi_1$ of type ${}^2B_2$ generated by two simple roots.  For a simple root $\alpha$, take the standard representative $n_\alpha = x_\alpha(1)x_{-\alpha}(-1)x_\alpha(1)$.
These representatives satisfy $n_\alpha^2 = h_\alpha(-1)$ and the corresponding braid relations; their conjugate action by root groups has the form
$$\dot{w}x_\beta(t)\dot{w}^{-1}=x_{w\beta}(\varepsilon_{w,\beta}t),\quad \varepsilon_{w,\beta}\in\{1,-1\}.$$ 
Since $-1=1$ in $\overline{\mathbb{F}}_2$, 
$n_\alpha^2=1$ and $ \varepsilon_{w,\beta}=1$.
Take $n=\dot{w}=\dot{w}_0\dot{w}_1$ and $n^2=1.$
Since $w\neq 1$, we have $n\neq 1$, and thus $n$ is an involution. For the above equation, we have 
$n x_\beta(t)n^{-1}=x_\beta(t)$ for $\beta\in\Phi_1$.  
These root groups generate the derived group of $M$, so $[n,{}^2B_2(q^2)]=1$. On the other hand, \(w_1\) acts trivially on $C_{q^2-1}$, while \(w_0=-1\) acts as inversion on the  maximal torus. Then \(w=w_0w_1\) also acts as inversion on $C_{q^2-1}$.  Thus for any $z\in C_{q^2-1}$, we have 
$z^n=z^{-1}$.  Let $F=F_2^f\gamma$ the Steinberg endomorphism defining the Ree group, where $\gamma^2=F_2$ and \(\gamma\) is the exceptional automorphism. Then $\gamma(n)=n$, and so $F_2(n)=n$. Thus
$n\in S$ and $n^d=n$ for any $d\in D$, as claimed.

Note that the root groups of $\bf J$ consist of  $\Phi_1$ and $\Phi^{+}\setminus\Phi_1^+$, where the latter generates ${\bf U}$. Since \(w_1\) preserves \(\Phi\setminus\Phi_1^+\), while \(w_0\) maps positive roots to negative roots, $w(\Phi^+\setminus\Phi_1^+) =\Phi^-\setminus\Phi_1^-$. Meanwhile, \(w\) fixes \(\Phi_1\) pointwise. Consequently, ${\bf J}^n=\bf J^-$ and $\bf J\cap\bf J^-=\bf M$, where \(\bf J^-\) is the opposite parabolic subgroup with the same Levi component $\bf M$. Taking \(F\)-fixed points yields $J\cap J^n=M$. This fact is used in Mackey's formula blew.

The character $W$ is $D_\chi$-fixed by Lemma \ref{defect zero characters of $K$}. So for any $d\in D_\chi$, Equation \eqref{eq: act} gives $\theta^d\in\{\theta,\theta^{-1}\}$. Suppose that $\theta^d=\theta^{-1}$.
Let $d$ has odd order $r$ because $D_\chi$ has odd order. Then $\theta=\theta^{d^r}=\theta^{(-1)^r}=\theta^{-1}$, and this implies $\theta^2=1$. However, the order of $C_{q^2-1}$ is odd, which forces $\theta=1$, contradicting the assumption. So $\theta$ is $D_\chi$-fixed. Therefore, \(\theta\) can be trivially extended to \(C_{q^2-1}\rtimes D_\chi\) by  $\widehat\theta(zd)=\theta(z)$ for any $z\in C_{q^2-1}$ and $d\in D_\chi$. 

Note that $W$ is unipotent by Lemma \ref{defect zero characters of $K$}. By Lemma \ref{lem:unipotent},   $W$ has 
an $(\HH_p)_W$-invariant extension to ${}^2B_2(q^2)\rtimes D$. Then restrict it to \({}^2B_2(q^2)\rtimes D_\chi\) and denote the restriction by \(\widehat W\). Define 
\[
\widehat\rho(zsd)=\theta(z)\widehat W(sd),
\]
where $z\in C_{q^2-1},\ s\in {}^2B_2(q^2)$ and $d\in D_\chi$.
Then  $\widehat\rho\in\Irr(M\rtimes D_\chi)$ and it's an extension of \(\theta\otimes W\). Inflate it via
$(J\rtimes D_\chi)/U\cong M\rtimes D_\chi$ to obtain \(\widehat\lambda\in\Irr(J\rtimes D_\chi)\), and let
$\widehat\chi=\Ind_{J\rtimes D_\chi}^{S\rtimes D_\chi}(\widehat\lambda)$. Clearly, $\widehat\chi$ is irreducible.
Since \(S\rtimes D_\chi=S(J\rtimes D_\chi)\) and \(S\cap (J\rtimes D_\chi)=J\), restriction yields
\[
\widehat\chi|_S=\Ind_J^S\Inf_M^J(\theta\otimes W)=\chi.
\]

Finally, we verify the stability. For any
$a=d\sigma\in(D\times\HH_p)_\chi$. Then by Equation \eqref{eq:hc},  $\theta^{d\sigma}\in\{\theta,\theta^{-1}\}$
and $W^{d\sigma}=W$. Since \(W^d=W\) by Lemma \ref{defect zero characters of $K$}, we have \(W^\sigma=W\). Therefore, the chosen extension satisfies $\widehat W^{d\sigma}=\widehat W$. If $\theta^{d\sigma}=\theta$, 
then by $\widehat \rho$ immediately yields
$\widehat\lambda^{\,a}=\widehat\lambda$. The induction commutes with field automorphisms and  Galois automorphisms, 
so $\widehat\chi^{\,a}=\widehat\chi$. If
$\theta^{d\sigma}=\theta^{-1}$. Since \(n\) centralizes \({ }^2B_2(q^2)\) and is $D$-fixed, it also centralizes \({}^2B_2(q^2)\rtimes D_\chi\). Hence, on the group \(M\rtimes D_\chi\), we obtain
\[
\widehat\rho^{\,a}(zsd)=
\theta^{-1}(z)\widehat W(sd)=
\widehat\rho^{\,n}(zsd).
\]
Moreover, since $n^d=n$ and $J\cap J^n=M$, 
$(J\rtimes D_\chi)\cap(J\rtimes D_\chi)^{\,n}=
M\rtimes D_\chi$. Apply Mackey's formula, we have
\[
\begin{aligned}
\left\langle
\widehat\chi^{\,a},\widehat\chi
\right\rangle_{S\rtimes D_\chi} &=
\sum_{x\in J\rtimes D_\chi\backslash S\rtimes D_\chi/J\rtimes D_\chi}
\left\langle
\widehat\lambda^{\,a}|_{(J\rtimes D_\chi)\cap (J\rtimes D_\chi)^x},\widehat\lambda^{\,x}|_{(J\rtimes D_\chi)\cap (J\rtimes D_\chi)^x}
\right\rangle,
\end{aligned}
\]
where the term corresponding to \(x=n\) equals
$\left\langle
\widehat\rho^{\,a},\widehat\rho^{\,n}
\right\rangle_{M\rtimes D_\chi}=
\left\langle
\widehat\rho^{\,n},\widehat\rho^{\,n}
\right\rangle_{M\rtimes D_\chi}=1$, and all other terms are non-negative integers. Thus 
$\left\langle
\widehat\chi^{\,a},\widehat\chi
\right\rangle_{S\rtimes D_\chi}\geq1$.
However, we know that $\widehat\chi^{\,a}$ and $\widehat\chi$
are irreducible, and this forces  $\widehat\chi^{\,a}=\widehat\chi$. We complete the proof.
\end{proof}

\begin{lemma}\label{lem: The non-split Suzuki-centralizer families}
If $\chi$ is any of characters in the families $\chi_{43},\chi_{44},\chi_{47},\chi_{48}$, then it exists an $(\HH_p\times D)_\chi$-invariant extension  to $S\rtimes D_\chi$.
\end{lemma}
\begin{proof}
Take the torus $T_\pm\cong C_{q^2\pm\sqrt{2}q+1}\times C_{q^2\pm\sqrt{2}q+1}$. In the notation of \cite[Table~B.2]{HH}, these are $T_6$ and $T_7$, and the torus $C_{q^2+\sqrt{2}q+1}\times C_{q^2-\sqrt{2}q+1}$ is $T_5$ with order $q^4+1$. For a non-trivial $\theta\in\Irr(T_\pm)$, the formulas in \cite[Appendix~B, construction of the $g_8$ and $g_{10}$ series]{HH} give
\begin{align}
 R_{T_-}^S(\theta)&=-\chi_{42}+\chi_{43}+\chi_{44}+\chi_{45},\label{eq:dlminus}\\
 R_{T_+}^S(\theta)&=-\chi_{46}+\chi_{47}+\chi_{48}+\chi_{49}.\label{eq:dlplus}
\end{align}
For convenience, from left to right, we will denote these four constituents as $\chi_1,\chi_2,\chi_3,\chi_4$.

First, we claim that $T_\pm$ is stable under a representative element $F_0$ of $D$ and that all rational Weyl representative elements can be taken $D$-fixed.
Let $\gamma$ be the exceptional endomorphism with $\gamma^2=F_2$. In the twisted case, write $F'=c_n\gamma^{2m+1}$ and $v=c_n\bar\gamma$, where $\bar{\gamma}$ is a finite-order action on the root system after removing the Frobenius scaling and $c_n$ 
denotes the inner automorphism action of $n$. Here, $n$ is
a preimage of an element of the Weyl group and $v$ has order $8$. Let $k=2m+1$. Since \(k\) is odd, there exists an integer \(b\) such that $bk\equiv1\pmod8$.
Hence, $(v^b)^k=v^{bk}=v$. Also, since \(b\) is odd, we may write $v^b=n_0\bar\gamma$, and define correspondingly
$F_0=c_{n_0}\gamma$. Then in characteristic $2$, 
the Tits kernel $\langle h_\alpha(-1)\rangle$ is trivial, so
we deduce that $F_0^k=F'$ and $F_0$ fixes all chosen  rational Weyl representative elements. That is, $F_0$ fixes $W^{F'}=\CC_{W}(v)$ pointwise, where $W=\NN_S(T_{\pm})/T_\pm$ is the rational Weyl group of $S$,
as claimed.

Next, for every $\chi$ lying in the families $\chi_i(i=1,2,3,4)$, through the following paragraphs, we  will construct it's extension $\widehat \chi\in\Irr(S\rtimes D_\chi)$.  

Write $H=W_\theta$ for the stabilizer of $\theta$ in $W$. Note that $\theta$ is trivial on $Z_{p'}=(C_{ q^2\pm\sqrt{2}q+1})_{p'}$,  so by the above discussion of \eqref{eq:rankone} , we have $H\cong C_4$. Also, by Table \ref{tab:1}, $W\cong G_8$ which has order $96$. Then $\ZZ(W)\cong C_4$. By \cite[Theorem 5.5]{MT18}, $\NN_W(H)=H\times \ZZ(W)$, so $\NN_W(H)/H\cong C_4$. Let \(d\in D_\chi\). Then \(d\) stabilizes the corresponding $S$-class of a semisimple element, and so we may write $\theta^d=\theta^{w_d}$ for some $w_d\in W$.
Thus  $W_{\theta^d}=W_\theta=H$ and \(w_d\) normalizes \(H\). This yields a group homomorphism $D\longrightarrow\NN_W(H)/H$, and the order of its image divides both \(|D_\chi|\) and \(4\). However, \(|D_\chi|\) is odd, so this image must be trivial. Consequently, for 
any $d\in D_\chi$, we have $\theta^d=\theta$. We now define
\[
\widehat\theta(td)=\theta(t),
\]
where $t\in T_{\pm}$ and $ d\in D_\chi$. Clearly, it's a linear character of $T_{\pm}\rtimes D_\chi$ and is 
an extension of $\theta$.

For every family \(\chi_i\) in Equations \eqref{eq:dlminus} and \eqref{eq:dlplus}, we claim that $\chi_i$ is $D$-fixed.  By the degrees in \cite[Table~A.2]{Himstedt},  we can partition them into three parts $\{\chi_1\}, \{\chi_2,\chi_3\}, \{\chi_4\}$, where the middle pair have the same degree, other parts have distinct degrees. So $\chi_i^d=\chi_i$ for any $d\in D$ and $i=2,4$. Also, since $D$ has odd order, it must act trivially on $\{\chi_2,\chi_3\}$. Thus every $\chi_i$ is $D$-fixed individually, as claimed. 

Since the action of $D$ fixes the chosen Weyl representative elements, every element of \(\NN_S(T_{\pm})\) can be written as $tn_w$, where $t\in T_{\pm}$ and $w\in W$ with each \(n_w\) commuting with \(D\). Therefore
$T_{\pm}\rtimes D\unlhd \NN_S(T_\pm)\rtimes D$ 
and  $\NN_S(T_\pm)\rtimes D/T_{\pm}\rtimes D\cong W$.
By Frobenius reciprocity and the restriction formula for normal subgroups,
\begin{equation}\label{2}
\begin{aligned}
\left\langle
\Ind_{T_\pm\rtimes D_\chi}^{\NN_S(T_\pm)\rtimes D_\chi}(\widehat \theta),
\Ind_{T_\pm\rtimes D_\chi}^{\NN_S(T_\pm)\rtimes D_\chi}(\widehat \theta)
\right\rangle_{\NN_S(T_\pm)\rtimes D_\chi}
&=
\sum_{w\in W}
\left\langle
\widehat\theta,\widehat\theta^{\,n_w}
\right\rangle_{T_\pm\rtimes D_\chi}.
\end{aligned}
\end{equation}
Since \(\widehat\theta\) is linear, each summand is either \(0\) or \(1\). Furthermore, since \(n_w\) commutes with \(D\), and  both sides take the value \(1\) on \(D_\chi\), we have that $\widehat\theta^{\,n_w}=\widehat\theta$ if and only if  $\theta^w=\theta$. Thus the Equation \eqref{2} is equal to $|H|=4$.

We apply some results of \cite[Sections 2-3]{DM} to prove  $\langle R_{T_\pm\rtimes D_\chi}^{S\rtimes D_\chi}(\widehat \theta), R_{T_\pm\rtimes D_\chi}^{S\rtimes D_\chi}(\widehat \theta) \rangle=4$.
Let $|D_\chi|=e$, $r=\frac{2m+1}{e}$ and $F_1=F_0^r$.
Then $F_1^e=F'$. Set ${\bf X}^\circ={\bf S}^e$, and let $\tau$ cyclically permute these $e$ factors
$\tau(x_1,\dots,x_e)=(x_e,x_1,\dots,x_{e-1})$.
Define 
\[
{\cal F}(x_1,\dots,x_e)=(F_1(x_e),F_1(x_1),\dots,F_1(x_{e-1}))
\]
Then
$x\mapsto (x,F_1(x),\dots,F_1^{e-1}(x))$
gives 
\begin{equation}\label{3}
S={\bf S}^{F'}\cong ({\bf X}^\circ)^{\cal F}.   
\end{equation}
Under this correspondence, the action of $\tau$ corresponds to the action of $F_1^{-1}$. Therefore, let
${\bf X}={\bf X}^\circ\rtimes\langle\tau\rangle$ and let ${\cal F}$ fix $\tau$. We then obtain the finite group structure 
\[
S\rtimes D_\chi\cong {\bf X}^F.
\]
Take a suitable $F_1$-stable maximal torus ${\bf T}$ in ${\bf S}$ contained in the $F_1$-stable Borel subgroup ${\bf B}$. 
Then the disconnected torus ${\bf L}={\bf T}^e\rtimes \langle \tau\rangle$ is well-defined and $\cal F$-stable.
Thus $T_{\pm}\rtimes D_\chi\cong{\bf L}^{\cal F}$ by the isomorphism $T_{\pm}\cong({\bf T}^e)^{\cal F}$.
As in \cite[Definition 2.4]{DM}, we choose 
$\widehat t\in\Irr({\bf L}^{\cal F})$ corresponding to $\hat{\theta}\in\Irr(T_{\pm}\rtimes D_\chi)$, and then construct the disconnected Lusztig induction
\[
\widehat R:=R_{{\bf L}^{\cal F}}^{{\bf X}^{\cal F}}(\widehat t)\in{\mathbb Z}\Irr({\bf L}^{\cal F}),
\]
and it is not difficult to verify that 
\begin{equation}\label{4}
{\rm Res}^{{\bf X}^{\cal F}}_{({\bf X^\circ})^{\cal F}}(\widehat R)=R_{({\bf T}^e)^{\cal F}}^{({{\bf X}^\circ})^{\cal F}}(t)=
R_{T_\pm}^S(\theta)=-\chi_1+\chi_2+\chi_3+\chi_4, 
\end{equation}
where $t\in\Irr(({\bf T}^e)^{\cal F})$ corresponds to  $\theta\in\Irr(T_\pm)$ by \eqref{3}. Moreover, by definition of ${\cal F}$, we know that it is also a Steinberg map. Apply Equation \eqref{2} and \cite[Proposition 3.2]{DM}, we have 
\begin{equation}\label{5}
\langle R_{T_\pm\rtimes  D_\chi}^{S\rtimes  D_\chi}(\widehat \theta), R_{T_\pm\rtimes  D_\chi}^{S\rtimes  D_\chi}(\widehat \theta) \rangle=\langle R_{{\bf L}^{\cal F}}^{{\bf X}^{\cal F}}(\widehat t), R_{{\bf L}^{\cal F}}^{{\bf X}^{\cal F}}(\widehat t) \rangle=\langle \widehat R, \widehat R \rangle_{{\bf X}^{\cal F}}=4.
\end{equation}

Since $D_\chi$ is cyclic and every $\chi_i$ is $D_\chi$-fixed, $\chi_i$ can be extended to $S\rtimes D_\chi$. Choose an extension $\chi_i^*$ and  for any  $\lambda\in\Irr(D_\chi)$, all $\chi_i^*\lambda$ are extensions of $\chi_i$. Therefore, we may write
\[
\widehat{R}=\sum_{i=1}^4\sum_{\lambda\in\Irr(D_\chi)}a_{i,\lambda}\chi_i^*\lambda+\Psi
\]
with $a_{i,\lambda}\in\mathbb{Z}$, and every irreducible constituent of $\Psi$ does not lie above any of the four $\chi_i$.
Restricting to $S$ and comparing the coefficients of $\chi_i$, we obtain
$\sum_{\lambda}a_{i,\lambda}=\varepsilon_i=\pm 1$.
Note that for each $i$, $\sum_{\lambda}a_{i,\lambda}^2\ge 1$.
On the other hand, together with Equation \eqref{5}, we have
\[
4=\langle\widehat{R},\widehat{R}\rangle_{S\rtimes D_\chi}
=\sum_{i=1}^4\sum_{\lambda}a_{i,\lambda}^2+\langle \Psi,\Psi\rangle_{S\rtimes D_\chi}\ge 4.
\]
This forces $\Psi=0$ and $\sum_{\lambda}a_{i,\lambda}^2=1$ for every $i$. Thus we obtain 
$\widehat{R}=\sum_{i=1}^4\varepsilon_i\widehat{\chi}_i$,
$\widehat{\chi}_i\big|_S=\chi_i$ and every $\widehat{\chi}_i$ is the unique irreducible character in $\widehat{R}$ lying above $\chi_i$.

Finally, we verify the stability. Let
$a=d\sigma\in (D\times \HH_p)_\chi$. 
Then $a$ fixes the corresponding Lusztig series.
So there exists some $w\in W$ such that $\theta^a=\theta^w$.
Since $w$ admits a $D$-fixed representative element $n_w\in\NN_S({\bf T})$, and $\widehat{\theta}$ takes value $1$ on $D_\chi$, we have
$\widehat{\theta}^a=\widehat{\theta}^{n_w}$.  Two compatibilities for Deligne–Lusztig induction ${\widehat R}$ on the above disconnected group $({\bf X}, \cal F)$ are required
\begin{equation}\label{7}
R_{{\bf L}^{\cal F}}^{{\bf X}^{\cal F}}(\widehat t)^a=R_{{\bf L}^{\cal F}}^{{\bf X}^{\cal F}}(\widehat{t}^a)
\ \text{and}\ R_{{\bf L}^{\cal F}}^{{\bf X}^{\cal F}}(\widehat{t}^{n})=R_{{\bf L}^{\cal F}}^{{\bf X}^{\cal F}}(\widehat{t}),
\end{equation}
where $n=(n_w,\ldots,n_w)\in(\mathbf X^\circ)^{\mathcal F}$ 
corresponds to $n_w$. By definition of ${R}_L^X$, and together with \cite[Proposition 3.3]{DL}, the first equation holds. Since $n_w\in\NN_S({\bf T})$ and $F_0(n_w)=n_w$,
we have 
$n{\bf T}^en^{-1}={\bf T}^e$ and $n\tau=\tau n$, so that
$n\mathbf L n^{-1}=\mathbf L$. 
Consider the inner automorphism $c_n:\mathbf X\rightarrow\mathbf X, x\mapsto nxn^{-1}$.
Since \(n\in X\), it commutes with \(\cal F\) and induces an isomorphism of the corresponding algebraic varieties.
Using the independence of the quasi-parabolic choice(see \cite[Theorem 4.8.13]{GM20}), and since $R_{{\bf L}^{\cal F}}^{{\bf X}^{\cal F}}(\widehat t)$ is a  class function, 
$R_{{\bf L}^{\cal F}}^{{\bf X}^{\cal F}}(\widehat{t}^{n})=R_{{\bf L}^{\cal F}}^{{\bf X}^{\cal F}}(\widehat{t})^n=R_{{\bf L}^{\cal F}}^{{\bf X}^{\cal F}}(\widehat{t})$.
Return to the group $S\rtimes D_\chi$, by Equation \eqref{7} and
$\widehat{\theta}^a=\widehat{\theta}^{n_w}$, we obtain
\begin{equation}\label{8}
R_{T_\pm\rtimes D_\chi}^{S\rtimes D_\chi}(\widehat \theta)^a=R_{T_\pm\rtimes D_\chi}^{S\rtimes D_\chi}(\widehat \theta^a)=R_{T_\pm\rtimes D_\chi}^{S\rtimes D_\chi}(\widehat \theta^{n_w})=R_{T_\pm\rtimes D_\chi}^{S\rtimes D_\chi}(\widehat \theta).    
\end{equation}
Also, since \(\widehat\chi\) is the unique irreducible character in \(\widehat R\) lying over \(\chi\) and
$(\widehat\chi^a)|_S=(\widehat\chi|_S)^a=\chi^a=\chi$.
Meanwhile, Equation \eqref{8} guarantees that \(\widehat\chi^a\) still appears in \(\widehat R\) with the same coefficient. Therefore, the uniqueness forces $\widehat\chi^a=\widehat\chi$. We complete the proof.
\end{proof}

All extensions of the weight characters are handled by the following proposition. 
\begin{proposition}\label{equ:local extention p>=5}
Let $\chi\in\IBr_p(S)$ and  $\Omega(\chi)=\overline{(R,\varphi)}\in\WW_p(S)$.
There exists $D_0\leq D$ such that $\varphi^\circ\in\IBr(\NN_S(R))$  has an $(\HH_p\times \NN_{S\rtimes D_0}(R))_\varphi$-invariant extension to $\NN_{S\rtimes D_0}(R)_{\varphi}$. 
\end{proposition}
\begin{proof}
Suppose that $R$ is non-cyclic abelian. Then by Table \ref{tab:1}, $\NN_S(R)/R=T_{p'}\rtimes W$, where $W$ is isomorphic to  $D_{16}$, $\GL_2(3)$, or $G_8$. Note that $p\geq 5$  and $p\nmid |\NN_S(R)/R|$, so $p\nmid\varphi(1)$. Thus we use \cite[Proposition~4.5]{Johansson},
$\varphi$ has an $(\HH_p\times \NN_{S\rtimes D}(R))_\varphi$-invariant extension $\widehat\varphi\in\Irr(\NN_{S\rtimes D}(R)_\varphi)$. The restriction $\widehat \varphi^\circ$ is
an $(\HH_p\times \NN_{S\rtimes D}(R))_\varphi$-invariant extension of $\varphi^\circ$.

Suppose that $R$ is cyclic. For convenience, we write $N=(Z_{p'}\rtimes E)\times K$ and $I=(Z_{p'}\rtimes E_\theta)\times K$. Then Equation \eqref{eq:rankone} gives 
\[
\varphi:=\psi_{\theta,\kappa,\epsilon}=
 \Ind_{I}^{N}
((\widehat\theta\epsilon)\otimes\kappa)\in\dz_p(\NN_S(R)/R),
\]
where $\widehat \theta\in\Irr(Z_{p'}\rtimes E)$ is
the trivial extension of $\theta\in\Irr(Z_{p'})$, $\epsilon\in\Irr(E_\theta)$ and $\kappa\in\dz_p(K)$.
Moreover, in the case, $\NN_S(R)=(Z\rtimes E)\times K$ and
$[E,K]=1$, where $Z$ and $K$ are as Table \ref{tab:2}, and $E=C_2$ or $C_4$.  By Clifford Theorem, we obtain $\operatorname{Res}_{Z_{p'}\times K}^{N}(\varphi)=\sum_{e\in E/E_\theta}\theta^e\otimes\kappa$.

Let $D_0\leq D$ be as in Lemma \ref{lem:3.5}. That is, $D_0$ has odd order, $D_0$ stabilizes $Z_{p'}$ and $K$, and satisfies $[D_0, E]=1$. For any $d\in (D_0)_\varphi$, $\kappa^d=\kappa$ and $\theta^d=\theta^e$ for some $e\in E$. Let the order of $d$ be the odd integer $m$. Since $[d,E]=1$, the action of field automorphism commutes with the $E$-action. This yields $\theta=\theta^{d^m}=\theta^{e^m}$. Therefore,
$\theta^d=\theta$ and $\kappa^d=\kappa$ for any $d\in (D_0)_\varphi$.

We construct an $(\HH_p\times D_0)_\kappa$-invariant extension $\widehat \kappa$ of $\kappa$ as follows.  By Table \ref{tab:2}, $K=\SL_2(q^2)$ or ${}^2B_2(q^2)$. For the former, it's all ordinary characters are real-valued. 
For the latter, by \cite[p.71 and Table 4.2.1]{S16}, it's $p$-defect zero characters are real-valued expect for two unipotent characters $W^{+}$ and $W^{-}$ when $p|q^2-1$.
Suppose that $\kappa$ is real-valued. Note that $(K\rtimes (D_0)_\varphi)/K\simeq (D_0)_\varphi$ is a cyclic group of odd order. Then by Lemma \ref{lem:real}, $\kappa$ has the unique real-valued extension $\widehat{\kappa}\in\Irr(K\rtimes (D_0)_\varphi)$. For any $a=\sigma d\in(\HH_p\times D_0)_\kappa$, \(\widehat\kappa^{a}\) is still a real-valued extension of \(\kappa\). Thus the uniqueness yields
$\widehat\kappa^{a}=\widehat\kappa$. Suppose that $K={}^2B_2(q^2)$ and $\kappa=W^{\pm}$. In this case, by Lemma \ref{lem:unipotent}, we can choose $\widetilde\kappa\in\Irr(K\rtimes D_0)$ such that it is an \((\HH_p)_\kappa\)-invariant extension of \(\kappa\). Let $\widehat\kappa=\operatorname{Res}^{K\rtimes D_0}_{K\rtimes  (D_0)_\varphi}(\widetilde\kappa)$. Then it is also an extension of \(\kappa\). Since $\kappa^d=\kappa$,
$\kappa^{a}=\kappa$ implies \(\kappa^\sigma=\kappa\).
On the other hand, the $d$-conjugation in \(K\rtimes D_0\) is an inner automorphism, and thus $d$ fixes \(\widetilde\kappa\). Consequently, we have also $\widehat\kappa^{a}=\widehat\kappa$. Therefore, in both cases, we obtain \(\widehat\kappa\) with the required simultaneous invariance.

We next construct an extension $\widehat \varphi$ of $\varphi$ as follows. Let $\widehat I=I\rtimes (D_0)_\varphi$ and $\widehat N=N\rtimes (D_0)_\varphi$. Since $\theta$ is $(D_0)_\varphi$-fixed and $[(D_0)_\varphi,E]=1$, we may write $\widehat I=(Z_{p'}\times K)\rtimes (E_\theta \times (D_0)_\varphi)$. 
We define  $\widetilde{\theta}(ted)=\theta(t)$ and 
$\widehat{\epsilon}(ed)=\epsilon(e)$, where 
$t\in Z_{p'}$, $e\in E_\theta$ and $d\in (D_0)_\varphi$.
Then $\widetilde \theta$ and $\widehat \epsilon$ are linear characters of $Z_{p'}\rtimes(E_\theta\times (D_0)_\varphi)$ and  $E_\theta\times (D_0)_\varphi$, respectively.
Define on $\widehat I$ 
\[
\widehat\beta(tked)=\theta(t)\epsilon(e)\widehat\kappa(kd),
\]
where $t\in Z_{p'}$, $k\in K$, $e\in E_\theta$, $d\in (D_0)_\varphi$.
Then  $\widehat\beta$ is a character of $\widehat I$.
Write $\beta=\widehat \theta\epsilon\otimes \kappa$,
and then $\operatorname{Res}_{I}^
{\widehat I}(\widehat\beta)=\beta$. Since $\beta$ is irreducible, $\widehat\beta$ is also irreducible. Now define
\[
\widehat\varphi=\operatorname{Ind}_{\widehat I}^{\widehat N}(\widehat\beta).
\]
Then $\widehat\varphi$ is a character of $\widehat N$.
Since $\widehat N=\widehat IN$ and $N\cap \widehat I=I$, 
the Mackey's formula gives
\[
\operatorname{Res}_{N}^{\widehat N}(\widehat\varphi)=\Ind_{I}^{N}(\operatorname{Res}_{I}^
{\widehat I}(\widehat\beta))=\Ind_{I}^{N}(\beta)=\varphi.
\]
Since $\varphi$ is irreducible, so $\widehat \varphi$  is indeed an extension of $\varphi$.

Finally, we verify the stability of $\widehat \varphi$.
Let $a=d\sigma\in(D_0\times\HH_p)_\varphi$.
Since $\varphi^a=\varphi$, the direct product decomposition and the uniqueness of Clifford correspondence, there exists some $e_a\in E$ such that
$$
\theta^a=\theta^{e_a},\ \kappa^a=\kappa,\ \epsilon^\sigma=\epsilon.
$$
Here, the last equality here can also be obtained ad follows.
Since  $\theta^\circ$ is the extension taking value 1 on the complement, it satisfies
$(\theta^\circ)^a=(\theta^a)^\circ$, $(\theta^\circ)^{e_a}=(\theta^{e_a})^\circ$.
Also, because $E$ is abelian, conjugation by $e_a$ acts trivially on $E_\theta$. Combined with the uniqueness of Clifford parameters and $\varphi^a=\varphi$, this gives
$\epsilon^\sigma=\epsilon$, as claimed. 

From the earlier choice of \(\widehat\kappa\), we have $\widehat\kappa^a=\widehat\kappa$.
Moreover, since $[E,(D_0)_\varphi]=1$ and $ [E,K]=1$, \(e_a\) normalizes \(I\) and acts trivially on the factor \(K\rtimes (D_0)_\varphi\). Therefore, by the definition of $\widehat \beta$, we obtain $\widehat\beta^a=\widehat\beta^{\,e_a}$.
By compatibility of induction with the action of $a$, 
\[
\widehat\varphi^a
=\Ind_{\widehat I}^{\widehat N}(\widehat\beta^a)
=\Ind_{\widehat I}^{\widehat N}(\widehat\beta^{e_a})
=(\Ind_{\widehat I}^{\widehat N}(\widehat\beta))^{e_a}
=\widehat\varphi^{e_a}
=\widehat \varphi,
\]
where the final equality holds because $e_a\in N \leq {\widehat N}$ and $\widehat \varphi$ is a class function on $\widehat N$. The restriction $\widehat \varphi^\circ$ gives the desired extension of $\varphi^\circ$.  We complete the proof.
\end{proof}

\begin{theorem}\label{thm:2F_4 for p>=5}
The simple groups ${}^2F_4(2^{2m+1})$ with $m\geq 1$  satisfy the inductive NAW condition at the prime $p\geq 5$.
\end{theorem}
\begin{proof}
The cyclic Sylow cases were settled in Lemma \ref{lem:3.2}. For $p\geq5$ with non-cyclic Sylow subgroups, by Theorem \ref{equ:bijection when p>=5}, we obtain an required equivalent bijection. By  \cite[Lemma~10.7(ii)]{FFZ}, Lemma \ref{lem:real}--\ref{lem: The non-split Suzuki-centralizer families} and together with the map \eqref{eq:basic}, we obtain the required invariant extensions of the global characters. All extensions of the local characters is obtained by Proposition \ref{equ:local extention p>=5}. Finally, we apply Proposition \ref{prop:criterion}. We are done.
\end{proof}


\subsubsection{Verification at the prime \texorpdfstring{$p=3$}{p=3}}\label{p=3}
%\ \ \quad Write \(v_p (n) \) for the $p$-adic assignment of $n$. That is, for prime number $p$, integer \(n \neq 0 \), \(v_p (n) \) is a non-negative integer $k$ that satisfies \(p^k \mid n, \ p^{k+1} \nmid n \). In particular, we write \(v_p (0)=\infty\). 

In this Subsection, we fix $p=3$. Let $a=v_3(q^2+1)$ and $b=(q^2+1)_{3'}$. That is, $q^2+1=3^a b$, where $(3,b)=1$ and $a\geq 1$.
Keep the groups $S$ and $D$ as above. We will use the block and radical-normalizer classification in \cite[Tables~3--4 and the following discussion]{Malle}, together with the decomposition matrix in \cite[Section 4]{Himstedt}.

\begin{lemma}\label{lem:h3}
Let $N\unlhd A$ with $A/N$ cyclic with odd order.
If $\theta\in\IBr_3(N)$ is $\HH_3$-invariant, then it has the unique $\HH_3$-invariant extension to $A_\theta$.
\end{lemma}
\begin{proof}
On the finite field of Brauer character values, let $\sigma\in\HH_3$ generate the action $\zeta\mapsto\zeta^3$. For an extension $\chi$, write $\chi^\sigma=\mu\chi$. A character $\lambda\chi$ is fixed precisely when $\lambda^3\mu=\lambda$.
Again the solution is unique in the odd-order group $X=\IBr_3(A/N)$. Uniqueness also gives invariance under automorphisms stabilizing the original character.
\end{proof}


\subparagraph{(i) The principal 3-blocks:}
Let $B_0$ be the principal block of $S$. We write $\IBr(B_0)$ for 
the set of Brauer characters $\phi\in\IBr_p(S)$ such that $\phi$ is contained in the block $B_0$. Similar sets that appear in the following text are defined accordingly.
It follows from  \cite[first row of Table~3]{Malle} that  $\IBr(B_0)$ has nine irreducible Brauer characters. The notation we used is the CHEVIE notation. By \cite[Section 4 and Table~C.5]{Himstedt}, we write 
\begin{equation}
\IBr_3(B_0)=\{\phi_1,\phi_4,\phi_{5,1},\phi_7,\phi_8,\phi_{10},\phi_{15},\phi_{18},\phi_{21}\},
\end{equation}
which corresponds to the ordinary character set, forms a unitriangular set, as follows.
\begin{equation}\label{principal 3-block}
\mathcal{B}_0=\{\chi_1,\chi_4,\chi_{5,1},\chi_7,\chi_8,\chi_{10},\chi_{15},\chi_{18},\chi_{21}\}.
\end{equation}

\begin{lemma}\label{eq:nine}
All Brauer characters in the block $B_0$ are rational-valued, and also fixed by every field automorphism.
\end{lemma}
\begin{proof}
For $i\in\{1,4,7,8,10,18,21\}$ in \eqref{principal 3-block}, every \(\chi_i\) is unipotent by \cite[Section 4]{Himstedt}. Moreover, it follows from \cite[Table A.1]{Himstedt} that the degree of $\chi_i$ uniquely determines it. On the other hand, Galois automorphisms preserve degrees of characters. This shows that \(\chi_i \) is rational-valued. 
For the non-unipotent character $\chi_{5,1}$(see \cite[Table A.2]{Himstedt}), by \cite[Table~C.5]{Himstedt}, it's easy to 
check that 
\[
\phi_{5,1}
=\chi_9^\circ-\phi_1+\phi_4+\phi_8-\phi_{10}.
\]
Note that by \cite[Table A.1]{Himstedt}, all terms on the right-hand side take rational values , so \(\phi_{5,1}\) is rational-valued.
For $\chi_{15}$, by \cite[Table~A.1]{Himstedt}, we have
$\QQ (\chi_{15})=\QQ(\sqrt{3})$. So $\chi_{15}$ is not rational-valued. However, $\chi_{15}^{\circ}=\chi_{16}^\circ$ is rational-valued. 
To see this, we take any $3'$-element \(g\in S\). Then $\chi_{15}\in\QQ(\zeta_{|g|})$, and thus $\chi_{15}(g)\in
\mathbb Q(\sqrt{3})\cap\mathbb Q(\zeta_{|g|})=\mathbb Q$. 

Therefore, the unitriangularity implies that the nine Brauer characters are rational-valued. Clearly, they are also fixed by every field automorphism.
\end{proof}


It follows from  \cite[Theorem 5.1]{Malle} that there are nine $3$-weights corresponding to nine $3$-Brauers in the block $B_0$. As in \cite[Table 4]{Malle} and \cite[(2B)]{An}, non-trivial $3$-radical subgroups, corresponding normalizers and  $3$-defect  zero characters in Table \ref{tab:3}. Here,  $P_3$ denotes the Sylow $3$-subgroup of $S$, $Q_8$ is a quaternion group of order $8$, and ${-}.{-}$ denotes the expansion of two groups. The last row because $b = (q^2 + 1)_{3'}$.
\begin{table}[ht]
\centering
\caption{Non-trivial \(3\)-radical subgroups and  $3$-defect  zero characters}\label{tab:3}
\begin{tabular}{@{}llllrr@{}}
\toprule
$R$ & $\NN_S(R)/R$ & The degrees of $\dz_3(\NN_S(R)/R)$& condition\\
\midrule
$P_3$ & $C_2\times C_2$ & $1,1,1,1$& $a\geq 2$\\
$3^{1+2}_+$ & $\SL_2(3)$ & $3$& $a\geq 2$\\
$3^{1+2}_+$ & $\SL_2(3).C_2$ & $3,3$& $a\geq 2$\\
$3^{1+2}_+$ & $Q_8\rtimes C_2$ & $ 1,1,1,1,2,2,2$& $a=1$\\
$C_{3^a}\times C_{3^a}$ &$(C_b\times C_b).\GL_2(3)$& $3,3$& $a\geq 1$ \\
\bottomrule
\end{tabular}
\end{table}
\begin{lemma}\label{lem:3.15}
Let $R$ be a radical subgroup in  Table \ref{tab:3}.
Then all defect zero characters of $\NN_S(R)/R$ are
$\HH_3$-invariant, and also fixed by every field automorphism.
\end{lemma}
\begin{proof}
Suppose that  $a\geq2$. For $C_2\times C_2$, its four linear characters are rational-valued. For $\SL_2(3)$, its the unique degree-$3$ character is inflated from $A_4\cong \PSL_2(3)$. For the last
row when $a>1$, the two degree-$3$ characters are trivial on the normal subgroup \(C_b\times C_b\), and then they are inflated from the quotient group \(\mathrm{GL}_2(3)\).
Also, two degree-$3$ characters of $ \GL_2(3)$ are also  inflated from $S_4$.
For \(H=\SL_2(3).C_2\), write $K=\mathrm{SL}_2(3)$.
Then $H/\ZZ(K)\cong S_4$ or $A_4\times C_2$. Thus two degree-$3$ characters of $H$ are inflated from $S_4$ of $A_4\times C_2$. Note that all characters of $S_4$ and $A_4$ are rational-valued, so these characters are rational-valued.

Suppose that $a=1$. For the last row, similar to the above paragraph. For $Q_8\rtimes C_2$, we consider the group 
\[
SD_{16}=\langle x,y\mid x^8=y^2=1,\ yxy^{-1}=x^3\rangle,
\]
and its every element is uniquely written as \(x^i\) or \(x^iy\) for \(i\in{\mathbb Z}/8{\mathbb Z}\).
For any $\sigma\in\HH_3$ and $\theta\in \Irr(SD_{16})$, we have $\theta(g)^\sigma=\theta(g^{3^j})$ for $g\in\{x^i,x^iy\}$. However, note that $3^j\equiv1\ \text{or}\ 3\pmod8$, and that every element of $SD_{16}$ is conjugate to its cube. Therefore,  every character in $\Irr(SD_{16})$
are $\HH_p$-invariant.

For a $D$-stable radical subgroup $R$, $D$ induces automorphisms on $\NN_S(R)/R$ and permutes the associated weight characters. For $C_2\times C_2$, by proof of \cite{Malle}, its four linear characters are $D$-fixed. 
For \(\SL_2(3)\), its the unique degree-$3$ character is $D$-fixed. For \(\SL_2(3).2\) and \(\GL_2(3)\), each have two relevant degree-$3$ characters. Since $D$ has odd order, it acts trivially on a set of two elements, so they are $D$-fixed respectively. For \(SD_{16}\), $\Aut(SD_{16})$ is a \(2\)-group, so any automorphism of odd order is trivial. Thus the automorphism induced by \(D\) is trivial, and this show that the seven weight characters are fixed. We complete the proof.
\end{proof}



\subparagraph{(ii) The non-principal $3$-blocks:}
By \cite[Theorem  B]{Johansson}, there exists an $\HH_p$-equivalent bijection
\[
\mathcal{E}(S,s)\longrightarrow \mathcal{E}(\CC_S(s),1)
\]
That is, if $\chi\in\mathcal{E}(S,s)$ has 
Jordan decomposition $(s,\nu)$ for $s$ a semisimple $3'$-element of $S$ and $\nu\in\mathcal{E}(\CC_S(s),1)$ a unipotent character of $\CC_S(s)$, denoted by  $\chi_{s,\nu}$,  then $\chi^{\sigma_p}$ has Jordan  decomposition $(s^p, \nu^{\sigma_p})$. 
Let $\eta\in\in\IBr_3(G)$ correspond to $\chi_{s,\nu}$ under the bijection as in \eqref{eq:basic}. Then in the following text, we also call that $\eta$ has Jordan  decomposition $(s,\nu)$ and $\eta^{\sigma_p}$ has Jordan  decomposition $(s^p,\nu^{\sigma_p})$.


Let $B$ be a non-principal positive-defect block of $S$ and
$s$ a semisimple $3'$-element of $S$. Then by \cite[Table 3]{Malle}, the centralizer $\CC_S(s)$ is a torus, or a direct product of $\PSL_2(q^2)$ and a torus. 
Thus the ordinary basic set of $\Irr(B)\cap \mathcal{E}(S,s)$ has one character or two characters, with Jordan decomposition  characters $1$ or $\{1,\St \}$. We may write $\{\chi_{s,1}\}$ or $\{\chi_{s,1},\chi_{s,\St}\}$ for the associated  basic set, where the latter can be obtained by \cite[Section 4 and Tables~D.2 and D.5]{Himstedt}.



\begin{lemma}\label{lem:3.19}
All Brauer characters in the block $B$ are real-valued. Moreover, suppose that $\eta\in\IBr_3(B)$ has Jordan decomposition $(s,\nu)$. Then for any $\sigma\alpha\in\HH_3\times D$, where $\sigma$ maps $3'$-th root of unity $\zeta$ to $\zeta^{3^j}$, 
$\eta^{\sigma\alpha}$ has Jordan decomposition $(s^{3^{j}\alpha^{-1}},\nu)$.
\end{lemma}
\begin{proof}
Take an $F$-stable maximal torus $\mathbf T$ containing $s$. In the Weyl group of ${}^2F_4$, the longest element satisfies
$w_0=-1$, so it acts on the torus by inversion
$t\mapsto t^{-1}$ for any $t\in {\bf T}$. Since $-1$ is a central element and fixed by the action defining the Ree group, it also lies in the rational Weyl group $W=\NN_S(T)/T$ of this torus $T={\bf T}^F$. Using Lang's theorem, there exists  $n\in \NN_{\mathbf S}(\mathbf T)^F$ such that $nsn^{-1}=s^{-1}$.

By \cite[Theorem B]{Johansson},
we have  $\overline{\chi_{s,\nu}} \longleftrightarrow
(s^{-1},\overline{\nu})$. Here, $-$ denotes the complex conjugation, $\nu=1$ or $\mathrm{St}$. Then both are real-valued. Combining with
$(s^{-1},\overline{\nu})\sim(s,\nu)$,
and thus $\overline{\chi_{s,\nu}}=\chi_{s,\nu}$.
Also,  by the unitriangular basic set as in \eqref{eq:basic}, the corresponding Brauer characters are real-valued. For any $\sigma_3\in \HH_3$ mapping $\zeta$ to $\zeta^{3}$, we have
$\chi_{s,\nu}^{\sigma_p}=\chi_{s^{3},\nu}$.
Field automorphisms have the analogous parameter
description by \cite[Theorem 4.1]{Malle}.
We complete the proof.
\end{proof}

By \cite[Tables 3-4]{Malle}, non-trivial $3$-radical subgroups and associated normalizers for the block $B$ in Table \ref{tab:4} below. We will separately consider these three cases.
\begin{table}[ht]
\centering
\caption{Non-trivial \(3\)-radical subgroups and normalizers}\label{tab:4}
\begin{tabular}{@{}lllrr@{}}
\toprule
$R$ & $\NN_S(R)/R$&condition \\
\midrule
$C_{3^a}\times C_{3^a}$ &$(C_b\times C_b).\GL_2(3)$& $a\geq 1$\\
$C_{3^a}$ & $\SL_2(q^2)\times D_{2b}$& $a\geq 2$\\
$C_3$ & $\PSU_3(q^2).C_2$&$a\geq 1$\\
\bottomrule
\end{tabular}
\end{table}

Similar to Jordan decomposition of characters, if a character $\psi$ of a finite group is parametrized by an object $x$, then we denote the character by \(\psi_x\). Of course, we also allow for all the actions we are considering on $x$.

\begin{lemma}\label{lem:3.20}
Let $R$ be a radical subgroup in  Table \ref{tab:4}.
Then all defect zero characters of $\NN_S(R)/R$  are real-valued. Moreover, suppose that  $\psi_x\in\dz_3(\NN_S(R)/R)$ is parametrized by $x$. Then there exists $D_0\leq D$ such that for any $\sigma\alpha\in \HH_3\times D_0$ stabilizing the block $B$, $\psi_x^{\sigma\alpha}$ is parametrized by $x^{\sigma\alpha}$.
\end{lemma}
\begin{proof}
(i) For $R\cong C_{3^a}\times C_{3^a}$, we have $\NN_S(R)/R\cong (C_{b}\times C_{b})\rtimes\GL_2(3)$. Here, the action of 
forming a semi-direct product by $x^z=x^{-1}$ for any 
$z\in \ZZ(\GL_2(3))$. By \cite[TABLE I]{An}, we write $T=C_{3^ab}\times C_{3^ab}$ for a maximal torus of $S$ containing $R$, and then $E:=\NN_S(T)/T\cong \GL_2(3)$, $T_{3}=R$ and $T_{3'}=C_{b}\times C_{b}$.
Let $\theta\in\Irr(T_{3'})$. Since $\theta$ is $E_\theta$-stable, we define an extension $\widehat \theta\in\Irr(T_{3'}\rtimes E_\theta)$ of $\theta$ by $\widehat \theta(te)=\theta(t)$ for any $t\in T_{3'}$ and $e\in E_\theta$. As in \cite[Sections 3 and 5]{Malle}, the corresponding weight characters are written as
\begin{equation}
\psi_{\theta,\epsilon}
= \Ind_{T_{3'}\rtimes E_\theta}^{\NN_{S}(R)}
(\widehat\theta\,\epsilon),    
\end{equation}
where $\epsilon\in\Irr(E_\theta)$. Since
$\psi_{\theta,\epsilon}(1)=|E:E_\theta|\epsilon(1)$ and
$|\NN_S(R)|_3=|E|_3=3$, we have that $\psi_{\theta,\epsilon}\in\dz_3(\NN_S(R))$ if and only if
$\epsilon\in\dz_3(E_\theta)$. 

Suppose that $\theta=1$. Then $E_\theta=E$ and the weight character $\psi_{1, \epsilon}$ are lifts of $\epsilon\in\Irr(E_\theta)$. Its $3$-defect zero characters are exactly two characters of degree $3$, which come from $S_4$. So both of which are rational-valued.

Suppose that $\theta\neq 1$ and $\psi_{\theta,\epsilon}\in\dz_3(\NN_S(R))$. Then we claim that $E_\theta=1$ or $C_2$. Note that $\theta^z=\theta^{-1}$, where $z=-I_2\in\ZZ(\GL_2(3))$.
If \(z\in E_\theta\), then 
\(\theta=\theta^{-1}\), i.e., \(\theta^2=1\). However, the order of \(\theta\) is odd, which implies \(\theta=1\), a contradiction. Thus $z\notin E_\theta$.
Also, every subgroup of \(\GL_2(3)\) with order $4$ contains 
\(z\). So the order of a Sylow \(2\)-subgroup of \(E_\theta\) is at most \(2\). Since \(|E|=48\), we have
$|E_\theta|\in\{1,2,3,6\}$. If \(3\mid|E_\theta|\), then 
\(|E_\theta|\le6\). For any \(\epsilon\in\Irr(E_\theta)\),
$\epsilon(1)^2\le|E_\theta|\le6$, so \(\epsilon(1)\le2\), which makes it impossible that $\epsilon(1)_3=|E_\theta|_3=3$. Thus $E_\theta=1$ or $C_2$, as claimed.

When $E_\theta=1$, then $\epsilon=1$, and $\psi_{\theta,1}$ is a weight character. If $E_\theta\cong C_2$, then the induced construction gives two weight characters, denoted by $\psi_{\theta,1}$ and $\psi_{\theta, sgn}$, where $sgn$ is
the symbol character. For any $\sigma\in\HH_3$, note that $\widehat\theta^\sigma=\widehat{\theta^\sigma}$, and that $1$, $sgn$ are $\sigma$-invariant. So $\psi_{\theta,\epsilon}^\sigma=\psi_{\theta^\sigma,\epsilon}$.
In particular, since $\overline{\psi_{\theta,\epsilon}}=\psi_{\theta^{-1},\epsilon}$, \(z\) centralizes \(E_\theta\) and inverts on \(R\), we have $(\widehat\theta\epsilon)^z=\widehat{\theta^{-1}}\epsilon$. 
Thus  $\psi_{\theta^{-1},\epsilon}=\psi_{\theta,\epsilon}$. 
This gives that $\psi_{\theta,\epsilon}$ is real-valued. 

Note that $(b,|\GL_2(3)|)=1$, complements $E_\theta$ are $T$-conjugate, and then the induced construction is independent of the complement. Thus $\psi_{\theta,\epsilon}$ is equivalent for every pair $(\theta,\epsilon)$ under the action of every field automorphism. Here, it acts trivially on $\epsilon$.

(ii) For $R\cong C_{3^a}$ when $a\geq2$,
$\NN_S(R)/R\cong \SL_2(q^2)\times D_{2b}$. 
Since $3\nmid 2b$, every weight character is written as
\begin{equation}
\psi_{\kappa,\beta}=\kappa\otimes \beta,    
\end{equation}
where $\kappa\in\dz_3(\SL_2(q^2))$ and $\beta\in\Irr(D_{2b})$.
On the hand, note that $q^2$ is even, so $\PSL_2(q^2)=\SL_2(q^2)$. Since $3|q^2+1$, $\kappa$ can be induced from  a non-trivial character $\lambda$ of the non-split torus $C_{q^2+1}$. That is, $\kappa=\Ind_B^{\PSL_2(q^2)}\Inf_{C_{q^2+1}}^B(\lambda)$,
where $B$ is an upper triangular Borel subgroup containing $C_{q^2+1}$. We denote by $\kappa=\kappa_\lambda$.  View $\lambda$ as a parameter, by the involution of the Weyl group, so we write $\lambda\sim\lambda^{-1}$ for $\lambda$ and $\lambda^{-1}$ as the same parameter. 
Thus  $\kappa_\lambda^\sigma=\kappa_{\lambda^\sigma}$ for any $\sigma\in\HH_3$. In particular, $\overline{\kappa_\lambda}=\kappa_{\lambda^{-1}}=\kappa_{\lambda}$. 

On the other hand, since \(b\) is odd, $\Irr(D_{2b})$ consists of two linear characters $1$, $sgn$ and $(b-1)/2$ characters $\rho_j$ of degree \(2\), where $j\in {\mathbb Z}/b{\mathbb Z}$ and $j\sim-j$. When $\beta\in\{1,sgn\}$, we have that $\bar \beta=\beta$ and $\beta^\sigma=\beta$. Otherwise, $\beta$ can be induced from a non-trivial character $\lambda_j$ of  $C_{b}$, where $\lambda_j(x)=\zeta_b^j$. However, since $\lambda_j\sim\lambda_{-j}$, we have
$\rho_j^\sigma=\rho_{3^rj}$ for any 
$\sigma\in\HH_p$ mapping $3'$-th root of unity $\zeta$ to $\zeta^{3^r}$. In particular, $\bar{\lambda}_j=\lambda_{-j}=\lambda_j$ and then $\bar \rho_j=\rho_{-j}=\rho_j$. 
Therefore, $\psi_{\kappa,\beta}^\sigma=\psi_{\kappa^\sigma,\beta^\sigma}$
and 
$\overline{\psi_{\kappa,\beta}}=\psi_{\bar{\kappa},\bar{\beta}}=\psi_{\kappa,\beta}$. 

For the relevant field automorphism $\alpha$ mapping $x$ to $x^{2^k}$, since $\kappa_\lambda^\alpha=\kappa_{\lambda^\alpha}$ and
$\beta^\alpha=\beta$ for $\beta\in\{1,sgn\}$ and otherwise 
$\rho_j^\alpha=\rho_{(-2)^kj}=\rho_{2^kj}$ because $\rho_{-j}=\rho_j$, thus $\psi_{\kappa,\beta}^\alpha=\psi_{\kappa^\alpha,\beta^\alpha}$.


(iii) For $R\cong C_3$, $\NN_S(R)/R\cong U\rtimes\langle\tau\rangle$, where $U=\PSU_3(q^2)$ and $\tau$ is a graph automorphism of order $2$. 
Note that $(3, q^2+1)=3$, so the order $|U|=q^6(q^6+1)(q^4-1)/3$. Since \(3|q^2+1\), we have $v_3(q^2-1)=0$ and $v_3(q^4-q^2+1)=1$. Also, $q^4-1=(q^2+1)(q^2-1)$ and $q^6+1=(q^2+1)(q^4-q^2+1)$, so $v_3(q^4-1)=a$ and $v_3(q^6+1)=a+1$. Thus $v_3(|U|)=-1+(a+1)+a=2a$. This implies that every $\theta\in\dz_3(U)$ if and only if $v_3(\theta(1))=2a$.

As in \cite[Sections~2.1 and~3.1]{Orevkov}, $8$ families of ordinary character degrees of ${\rm GU}_3(q^2)$ are listed,
using their notations, as follows.
\[
\{1,q^2(q^2-1),q^6,q^4-q^2+1,\ q^2(q^4-q^2+1),\ (q^2-1)(q^4-q^2+1),\ d_7=q^6+1,\ d_8=(q^2-1)(q^2+1)^2\},
\]
and after restricting to \({\rm SU}_3(q^2)\), other two degrees generated $(q^2-1)(q^4-q^2+1)/3$ and $d_8/3$.
Therefore,  it's easily checked that the $3$-defect zero characters of $U$ come from the two families of degrees $d_7$ and $d_8$. For $d_7$-family, we have $v_3(d_7)=a+1$. Then $\theta\in\dz_3(U)$ if and only if $a=1$.  For $d_8$-family, we have $v_3(d_8)=2a$ and $v_3(d_8/3)=2a-1$. The other families, including the restriction generated degrees(see \cite[Section 3.1]{Orevkov}), none of them has a \(3\)-defect zero character.

Suppose that $a=1$. We will give the parametrization of characters $\psi_{u,z}$ of the $d_7$-family as follows.
To avoid confusion, we write \(d_7\)-family in \cite{Orevkov} as $X_{t,s}$, where $t\in{\mathbb Z}/(q^2+1)$, $s\in{\mathbb Z}/(q^4-1)$, $s\not\equiv0\pmod{q^2-1}$ and 
$X_{t,s}=X_{t,-q^2s}$. Multiplying by the linear character $\lambda$ obtained from the determinant changes the parametrization via
$(t,s)\longmapsto(t+1,s+1-q^2)$. After restriction to \({\rm SU}_3(q^2)\), $\lambda=1$ i. Thus we may normalize $s$ to \(0\) and retain  $k=s+(q^2-1)t\pmod{q^4-1}$. The equivalence relation after restriction becomes $k\sim -q^2k$, and thus
$X_{0,k}$ denotes characters of the \(d_7\)-family.

We now check whether we can descend to \(U\). 
Let $eI_3\in \ZZ({\rm SU}_3(q^2))$ be a  generator. Then it acts by \(e^k\) on the corresponding representation. So
$\ZZ({\rm SU}_3(q^2))\leq\ker(X_{0,k}|_{{\rm SU}_3(q^2)})$
if and only if $3\mid k$. We write $k=3j$, where 
$j\in{\mathbb Z}/\frac{q^4-1}{3}=\mathbb Z/((q^2-1)b)$.
Since $(q^2-1,b)=1$, the Chinese Remainder Theorem gives a bijection
\[ 
j\longleftrightarrow (u,z)\in {\mathbb Z}/(q^2-1)\times{\mathbb Z}/b.
\]
The condition \(k\not\equiv0\pmod{q^2-1}\) is equivalent to \(u\neq0\). Since $k\sim -q^2k$, $-q^2\equiv-1\pmod{q^2-1}$ and $-q^2\equiv1\pmod b$, we obtain $(u,z)\sim(-u,z)$. Therefore,
\begin{equation}
\{\psi_{u,z}|\ 0\ne u\in\mathbb Z/(q^2-1),\ z\in\mathbb Z/b,\ (u,z)\sim(-u,z)\}.    
\end{equation}

For $\tau$, which
can be realized via inverse transposition. So
$\psi_k^\tau=\psi{-k}=\overline{\psi_k}$.
Also  $-k\longleftrightarrow(-u,-z)\sim(u,-z)$, so
$\psi_{u,z}^{\tau}=\psi_{u,-z}=\overline{\psi_{u,z}}$.
For $\sigma\in \HH_3$ mapping any $3'$-th root of unity $\zeta$ to $\zeta^{3^r}$. Since \(a=1\), $v_3(q^4-1)=1$,
and then after substituting \(k=3j\), the order of the roots of unity involved in the eigenvalues divides
$\frac{q^4-1}{3}=(q^2-1)b$, which is coprime to \(3\).
Thus $\psi_k^\sigma=\psi_{3^rk}$, and this shows 
$\psi_{u,z}^{\sigma}=\psi_{3^ru,\,3^rz}$. 
In particular, $\overline{\psi_{u,z}}=\psi_{u,-z}=\psi_{u,z}$.
For  the relevant field automorphism $\alpha$ mapping $x$ to $x^{2^n}$, note that $\psi_k^\alpha=\psi_{(-2)^nk}=\psi_{2^nk}$, so $\psi_{u,z}^{\alpha}=\psi_{2^nu,\,2^nz}$.

For $d_8$-family, We give the parametrization of characters $\rho_v$ of the $d_8$-family as follows. Let $c=(q^4-q^2+1)/3$. Similar to the $d_7$-family, starting from the \(d_8\)-family of ${\rm GU}_3(q^2)$, whose parametrization is $t\sim q^4t$ and  $t\in{\mathbb Z}/(q^6+1)$.
Restriction to \({\rm SU}_3(q^2)\), the parametrization can be taken as $k\sim q^4k$ and $k\in{\mathbb Z}/(q^4-q^2+1)$.
Here, \(k=0\) does not yield the required irreducible characters, and the two special parametrization $k=c$ and $k=2c$ give  degree-\(d_8/3\) characters. As for the $d_7$-family, the central element \(eI_3\) acts by the scalar \(e^k\), so the condition for descent to \(U\) remains
$3\mid k$. However, \(3\nmid c\), neither of the two special 
parametrization \(c,2c\) satisfies this condition. 
We then write $k=3v$, where $v\in\mathbb Z/c$.
Excluding \(k=0\), we obtain 
\begin{equation}
 \{\psi_v|\ 0\ne v\in\mathbb Z/c,\ v\sim q^4v\}.   
\end{equation}
All these characters have degree
$\psi_v(1)=(q^2-1)(q^2+1)^2$ and are $3$-defect zero characters.
Similar to the $d_7$-family, for $\sigma$, $\tau$ and $\alpha$ as above, we have 
$\psi_{v}^{\tau}=\psi_{-v}=\overline{\psi_{v}}$,
$\psi_{v}^{\sigma}=\psi_{3^rv}$ and $\psi_{v}^{\alpha}=\psi_{2^nv}$.  
We complete the proof.
\end{proof}

\begin{theorem}\label{thm:2F4-p=3}
The simple groups ${}^2F_4(2^{2m+1})$ with $m\geq 1$  satisfy the inductive NAW condition at the prime $p=3$.
\end{theorem}
\begin{proof}
For the principal $3$-block $B_0$, \cite[Theorem 5.1]{Malle} gives
an bijection $\Omega_0$ between $\IBr_3(B_0)$ and the set of weights with the group $R$ in Table \ref{tab:3} as the radical subgroup. The equivalence of $\Omega_0$ under Galois automorphisms and field automorphisms is  verified from Lemmas \ref{eq:nine} and \ref{lem:3.15}. Apply Lemma \ref{lem:h3} to $A=S\rtimes D$ and $\NN_{S\rtimes D}(R)$, respectively, we obtain the required extensions condition \eqref{extension criterion} in global and local cases.

For the non-principal $3$-block $B$. Similarly, there is an bijection $\Omega_B$ between $\IBr_3(B)$ and the set of weights with the group $R$ in Table \ref{tab:4} as the radical subgroup.
Let $\eta\in\IBr_3(B)$ with the parametrization  $(s,\nu)$. Then by Lemma \ref{lem:3.19}, $\sigma\in\HH_3$ acts on $\eta$ by the action of $\sigma$ on $(s,\nu)$. Let $\overline{(R,\psi)}=\Omega_B(\eta)$. Here, $\psi$ is defined as in the proof of Lemma \ref{lem:3.20}. Then $\sigma$ acts on $\psi$ by acting on its the parametrization $x$. We can choose $x$ corresponding one to one with the global parametrization $(s,\nu)$.
By Lemmas \ref{lem:3.19} and \ref{lem:3.20}, $\Omega_B$ is
equivalent with respect to $\HH_3$ and  the set $D_0$ of some representatives of field automorphism in $D$. 

For $\eta$, which is real-valued by Lemmas \ref{lem:3.19},
So we apply Lemma \ref{lem:real} to $A=S\rtimes D_0$ and $N=S$, we obtain an $(\HH_3\times D_0)_\eta$-invariant extension $\widehat \eta\in\IBr(S\rtimes (D_0)_\eta)$ of $\eta$. For $\psi$, which is real-valued by Lemmas \ref{lem:3.20},
So we apply Lemma \ref{lem:real} to $A=\NN_{S\rtimes D_0}(R)$ and $N=\NN_S(R)$, we obtain an $(\HH_3\times \NN_{S\rtimes D_0}(R))_\psi$-invariant extension  $\widehat{\psi^\circ}\in\Irr(\NN_{S\rtimes D_0}(R)_\psi)$ of $\psi^\circ$.

Finally, together with the case of the defect-zero $3$-block, Proposition \ref{prop:criterion} holds. We complete the proof.
\end{proof}






\section{Some special groups}\label{section 4}
We now handle two exceptional cases ${}^2G_2(3)$ and ${}^2F_4(2)$, which are mentioned at the beginning of Section \ref{section 3}.

Let \(S={}^2F_4(2)'\) and \(L={}^2F_4(2)\). Then $L\cong \Aut(S)$ and we identify \(S\) with \(\Inn(S)\). Note that the order $|S|$ has the prime factors $p=2,3,5,13$ by \cite{T}. Moreover,  $\Out(S)=L/S\cong C_2$ which is only generated by an involution, and  $\CC_L(S)=1=\ZZ(L)$.
Since the Schur multiplier of $S$ is trivial, its own universal \(p'\)-covering group. We first will consider $p=3$ and $p=5$, respectively.

\subsection{Verification at the prime 
\texorpdfstring{$p=3$}{p=3}}\label{sec:p3}
Let $P=P_3\in\Syl_3(S)$. Then by \cite{T}, $P\cong3_{+}^{1+2}$ is the  extraspecial $3$-group of order $27$, $\ZZ(P) = P'\cong  C_3$ and $\ZZ(P)/P\cong {\mathbb F}_2^3$. Moreover,  $\NN_S(P)\cong P\rtimes D_8$ has order $216$ and $\NN_L(P)\cong P\rtimes SD_{16}$ has order $432$.  The quotient $SD_{16}$ embeds in $\GL_2(3)$ via the action on $P/\ZZ(P)$. Let
\begin{equation}\label{p3:eq:SDmat}
 u=\begin{pmatrix}1&1\\-1&1\end{pmatrix}\ \text{and}\
v=\begin{pmatrix}1&0\\0&-1\end{pmatrix}\ \text{over }\FF_3
\end{equation}
with \(u^8=v^2=1\) and \(vuv^{-1}=u^3\). 
Then \(D_8\cong\langle u^2,v\rangle\) and
\(SD_{16}\cong \langle u,v \rangle\).

\subparagraph{(i) Global characters and two actions.}
The 3-defect-zero ordinary characters of \(S\) are given by the ordinary characters table (see \cite{T}) and the \(3\)-decomposition matrix (see \cite[Block.1]{R}), we use their notation,
\begin{equation}\label{p3:eq:dzS}
 \dz_3(S)=\{\chi_4,\chi_5,\chi_9,\chi_{10},\chi_{11},
 \chi_{15},\chi_{16},\chi_{17},\chi_{20}\}.
\end{equation}
with degrees $27,27,351,351,351,675,702,702,1728$. 
The principal \(3\)-block $B_0$ of $S$ has nine Brauer characters, with labels and degrees(see \cite{R})
\begin{equation}\label{p3:eq:principal}
 \IBr_3(B_0)=\{\varphi_1,\varphi_2,\varphi_3,
 \varphi_6,\varphi_7,\varphi_8,\varphi_{12},
 \varphi_{16},\varphi_{17}\}
\end{equation}
with degrees $1,26,26,77,124,124,572,1099,1099$.
The other nine Brauer characters are exactly the restrictions \(\chi^\circ\) with \(\chi\in\dz_3(S)\) in \eqref{p3:eq:dzS}.   
%$|\IBr_3(S)|=18=|\mathcal W_3(S)|$.

Using the decomposition matrix in \cite{R} and the ordinary character values and class-power maps in \cite{T,OrdData}, one reconstructs all Brauer values on the eighteen \(3\)-regular classes. The third-power map $\pi_3$ is displayed in Appendix \ref{app:maps}. The fields generated by the values of the nine characters in the block $B_0$ are
\begin{table}[ht]
\centering
\caption{Field of Brauer characters in $B_0$}
\label{p3:tab:weights}
\begin{tabular}{@{}lllrr@{}}
\toprule
Brauer characters&Field of values\\
\midrule
\(\varphi_1,\varphi_6,\varphi_{12}\)&\(\QQ\)\\
\(\varphi_2,\varphi_3\)&\(\QQ(\sqrt{-2})\)\\
\(\varphi_7,\varphi_8,\varphi_{16},\varphi_{17}\)&\(\QQ(\sqrt{13})\)\\
\bottomrule
\end{tabular}
\end{table}
\begin{lemma}\label{equ}
All Brauer characters in the block $B_0$ are \(\mathcal H_3\)-invariant. Moreover, $\HH_3$ and $\Out(S)$ act respectively on $\IBr_3(S)$ by the non-trivial permutations
\[
\tau=(\varphi_4\ \varphi_5)(\varphi_{10}\ \varphi_{11})
 (\varphi_{14}\ \varphi_{15}),
\]
\[
\alpha=(\varphi_2\ \varphi_3)
 (\varphi_7\ \varphi_8)
 (\varphi_{14}\ \varphi_{15})
 (\varphi_{16}\ \varphi_{17}),
\]
and other characters are $(\HH_3\times \Out(S))$-invariant.
\end{lemma}
\begin{proof}
For $\sigma_3\in\HH_3$ mapping $\zeta$ to $\zeta^3$, where $\zeta$ is $3'$-th root of unity.
Since $\sqrt{-2}=\zeta_8+\zeta_8^3$, we obtain $\sigma_3(\sqrt{-2})=\zeta_8^3+\zeta_8^9=\sqrt{-2}$.
Let $\zeta_{13}$ is $13$-th root of unity.
Define $\chi(a)=\left(\frac{a}{13}\right)$,
where $\left(\frac{a}{13}\right)$ is the Legendre symbol. The corresponding quadratic Gauss sum is
$\tau_{13}=\sum_{a=1}^{12}\chi(a)\zeta_{13}^{a}$.
Since $13\equiv 1\pmod{4}$, the classical formula for quadratic Gauss sums gives $\tau_{13}=\sqrt{13}$.
Thus, $\tau_{13}$ generates 
$\mathbb{Q}(\tau_{13})=\mathbb{Q}(\sqrt{13})$.
Note that \(3\equiv4^2\pmod{13}\), so 
\(\sigma_3(\tau_{13})=\tau_{13}\). Thus we have
\(\sigma_3(\sqrt{13})=\sqrt{13}\).    
By \cite[Tables]{T}, the action of $\sigma_3$ on the set $\IBr_3(S)$ by the non-trivial permutation is $\tau=(\varphi_4\ \varphi_5) (\varphi_{10}\ \varphi_{11})(\varphi_{14}\ \varphi_{15})$, and  all other \(\varphi_i\) are $\HH_3$-fixed. Moreover, these three pairs occur in the $3$-defect-zero part with degrees \(27,351,702\), respectively.

Let \(1\neq \alpha\in \Out(S)\cong C_2\). By the stored class fusion of $S$ into $L$, which refers to a mapping from the conjugate class set of $S$ to the conjugate class set of $L$ and is recorded in \cite{OrdData} or Appendix \ref{app:maps}, and  together with two decomposition matrices of $S$ and $L$(see \cite{T,R}), we have the permutation
$\alpha=(\varphi_2\ \varphi_3) (\varphi_7\ \varphi_8)(\varphi_{14}\ \varphi_{15})(\varphi_{16}\ \varphi_{17})$. 
Thus in the block $B_0$ has three  $\alpha$-fixed points $\varphi_1,\varphi_6,\varphi_{12}$ and three $\alpha$-orbits
$\{\varphi_2\ \varphi_3\}, \{\varphi_7\ \varphi_8\}, \{\varphi_{16}\ \varphi_{17}\}$.
We complete the proof.
\end{proof} 



\subparagraph{(ii) Local characters and two actions.} By \cite[Theorem 5.1 and its proof]{Malle} and \eqref{p3:eq:dzS}, there are nine  $3$-weights with trivial radical group. These nine are separate $3$-defect-zero blocks.
Moreover, there are nine $3$-weights with non-trivial radical group $R$ as follows.
\begin{lemma}\label{p3:prop:radical}
There are exactly three non-trivial \(S\)-conjugacy classes of \(3\)-radical subgroups, represented by $R_+, R_-, P$, where \(R_\pm\cong C_3\times C_3\) satisfies
 $\NN_S(R_\pm)/R_\pm\cong\GL_2(3)$, and the non-trivial outer automorphism exchanges the two \(S\)-classes \([R_+]\) and \([R_-]\).
\end{lemma}
\begin{proof}
Since \(\NN_S(P)/P\cong D_8\) is a \(3'\)-group, \(\OO_3(\NN_S(P))=P\), and then $P$ is radical.
By \cite{A}, the ATLAS has a single class \(3X\) of elements of order \(3\), with \(|\CC_S(z)|=108\) for \(z\in3X\).
Since \(Z(P)\) is characteristic in \(P\),
\(\NN_S(P)\leq\NN_S(Z(P))\).  On the other hand,
\(\Nor_S(\langle z\rangle)/\Cen_S(z)\leq C_2\). Thus
\[
 216=|\NN_S(P)|\leq|\NN_S(Z(P))|\leq2\times108=216.
\]
Thus \(\NN_S(Z(P))=\NN_S(P)\), whose normal \(3\)-core is \(P>Z(P)\). All subgroups of order \(3\) are conjugate, so they are not radical.

The GAP maximal-subgroup list (see \cite{T}) contains \(L_3(3):2\leq S\). Inside its normal subgroup \(L_3(3)=\SL_3(3)\). Let \(Q=R_0:\GL_2(3)\) with \(R_0\cong C_3\times C_3\), and then \(|Q|=432\).  Take its Sylow \(3\)-subgroup \(P\) so that \(R_0\leq P\).  Since \(R\leq\NN_S(R_0)\) and  $|\NN_S(R_0)|\leq 432$, this forces
\(\NN_S(R_0)=Q\).  Thus \(R_0\) is radical because
\(\OO_3(\GL_2(3))=1\).

Note that the action of \(\NN_L(P)/P\cong SD_{16}\) on the four order-\(9\) subgroups of \(P\) is transitive by \eqref{p3:eq:SDmat}.  Thus they
are radical and have normalizer quotient \(\GL_2(3)\).
Note that \(\NN_S(P)/P=D_8\) has exactly two orbits on
the four subgroups.  These remain distinct under \(S\)-conjugation: if \(R^g=R'\) with \(R,R'\leq P\), then both \(P^g\) and \(P\) are Sylow \(3\)-subgroups of \(\NN_S(R')\).  Sylow conjugacy inside \(\NN_S(R')\) supplies \(h\in\NN_S(R')\) with \(P^{gh}=P\) and \(R^{gh}=R'\).  Therefore \(gh\in\NN_S(P)\), so the two subgroups belonged to the same \(D_8\)-orbit already. The transitivity of \(\NN_L(P)\) shows that the outer involution interchanges these two \(S\)-classes.
\end{proof}


\begin{lemma}\label{equ1}
Let $R$ be a radical subgroup in Lemma \ref{p3:prop:radical}.
Then all defect zero characters of $\NN_S(R)/R$  are rational-valued.
\end{lemma}
\begin{proof}
If \(R=R_\pm\), then \(\Nor_S(R)/R=\GL_2(3)\) by Lemma \ref{p3:prop:radical}, whose irreducible ordinary degrees are $1,1,2,2,2,3,3,4$ by \cite{T}.
Since \(|\GL_2(3)|_3=3\),  its two characters of degree
\(3\) are \(3\)-defect zero. They are inflated from the
two degree-\(3\) characters of \(\operatorname{PGL}_2(3)\cong S_4\). Thus they are rational-valued. If \(R=P\), then \(\Nor_S(R)/R=D_8\), which is a \(3'\)-group and has five ordinary irreducible characters with degrees \(1,1,1,1,2\) by \cite{T}. Moreover, they are $3$-defect zero characters and rational-valued. 
\end{proof}

Let \(1\neq \alpha\in \Out(S)\). For $R=P$, write
\(D_8=\langle r=u^2,s=v\rangle\) by using \eqref{p3:eq:SDmat}. Denote its linear characters by
\(\lambda_{\epsilon,\delta}(r)=\epsilon\),
\(\lambda_{\epsilon,\delta}(s)=\delta\) for
\(\epsilon,\delta\in\{\pm1\}\), and its degree-\(2\) character by \(\rho\).  
With the convention \(d^u=u^{-1}du\), conjugation sends
\(r\mapsto r\), \(s\mapsto rs\).  Thus, $\alpha$ fixes
 $\lambda_{+,+},\ \lambda_{+,-},\ \rho$ and interchanges \(\lambda_{-,+}\) with \(\lambda_{-,-}\).
For \(R=R_\pm\), we denote the two degree-\(3\) characters of \(\GL_2(3)\) by \(\tau_{\pm,0}\) and \(\tau_{\pm,1}\) such that $\alpha$ sends \(\tau_{+,j}\) to \(\tau_{-,j}\), \(j=0,1\). 
%Note that all characters of \(\GL_2(3)\) are rational-valued expect for two degree-\(2\) characters(their values contain in $\QQ(\sqrt{-2})$). Thus as the proof of Lemma \ref{equ}, all characters of \(\GL_2(3)\) are \(\HH_3\)-fixed.
Also, by Lemma \ref{p3:prop:radical}, $\alpha$ exchanges 
\(Q_+\)-weights and \(Q_-\)-weights. Consequently, the nine non-trivial weights have the same $\alpha$-orbit structure as \(\IBr_3(B_0)\), that is, we have the following correspondence.
\begin{table}[ht]
\centering
\caption{One correspondence on $B_0$}\label{tab:6}
\begin{tabular}{@{}ll@{}}
\toprule
$3$-Brauer characters&$3$-Weights\\
\midrule
\(\varphi_1\)&\(\overline{(P,\lambda_{+,+})}\)\\
\(\varphi_6\)&\(\overline{(P,\lambda_{+,-})}\)\\
\(\varphi_{12}\)&\(\overline{(P,\rho)}\)\\
\(\varphi_2,\ \varphi_3\)&\(\overline{(P,\lambda_{-,+})},\ \overline{(P,\lambda_{-,-})}\)\\
\(\varphi_7,\ \varphi_8\)&\(\overline{(Q_+,\tau_{+,0})},\ \overline{(Q_-,\tau_{-,0})}\)\\
\(\varphi_{16},\ \varphi_{17}\)&\(\overline{(Q_+,\tau_{+,1})},\ \overline{(Q_-,\tau_{-,1})}\)\\
\bottomrule
\end{tabular}
\end{table}

\begin{proposition}\label{p3:prop:equivariant}
There exists an 
\((\HH_3\times\Aut(S))\)-equivalent bijection 
\(\Omega:\IBr_3(S)\longrightarrow\mathcal W_3(S)\).
\end{proposition}
\begin{proof}
Define \(\Omega\) on the defect-zero blocks by
$\chi^\circ\mapsto\overline{(1,\chi)}$ for every $\chi\in\dz_3(S)$ in \eqref{p3:eq:dzS}. The map commutes with both Galois automorphisms and outer automorphisms. On the block \(B_0\), the action of $\HH_3$ is trivial on both sides by Lemmas \ref{equ} and \ref{equ1}. Moreover, by two permutations
in Lemma \ref{equ}, and together with actions on the weight characters, we  establish the required bijection.
\end{proof}

\subparagraph{(iii) Compatible extensions.}
The following verifications concern extensions of Brauer characters and weight characters for the above bijection. Choose \(\AT=L_{\eta^{\HH_3}}\) and \(A=L_\eta\) as in Proposition \ref{prop:criterion}.
\begin{proposition}\label{thm:p3}
Let $\eta\in\IBr_3(S)$ and $\Omega(\eta)=\overline{(R,\psi)}\in\WW_3(S)$.
Then $\eta$ and $\psi^\circ$ satisfy the condition \eqref{extension criterion}.
\end{proposition}
\begin{proof}
Note that $\eta$ comes from restriction of \eqref{p3:eq:dzS} and \eqref{p3:eq:principal}, and that $\psi$ comes from Table \ref{tab:6}. First, for the fixed three characters \(\varphi_1,\varphi_6,\varphi_{12}\) in the principal block $B_0$, their degrees are \(1,77,572\) by \eqref{p3:eq:principal}. Let \(\widehat\chi_6\) and \(\widehat\chi_{14}\) be ordinary rows \(8\) and \(21\) of \(L\) in \cite{OrdData}. They have degrees \(78\) and \(650\), and their restriction  to $S$ are \(\chi_6\) and \(\chi_{14}\). Moreover,  \(\widehat\chi_6\) and \(\widehat\chi_{14}\) are rational-valued. The subscripts here are their \(S\)-labels. The pertinent rows of \cite[Block. 1]{T} are
\[
 \chi_{6,0}^{\circ}=\varphi_{1,0}+\varphi_{6,0},\quad
\chi_{14,0}^{\circ}=\varphi_{1,0}+\varphi_{6,0}+\varphi_{12,0}.
\]
Here, the double subscripts are the labels for \(L\).
Choose the signs of the extensions compatibly with these rows. Then 
\[
\widehat\varphi_1=1_L,\ \widehat\varphi_6=\widehat\chi_6^{\circ}-1_L,\ 
\widehat\varphi_{12}=\widehat\chi_{14}^{\circ}-\widehat\chi_6^{\circ}
\]
are irreducible Brauer extensions of the required degrees. In particular, irreducibility of these
differences uses the decomposition matrix of $L$, not their degrees alone. All three are rational-valued, and thus
\(\HH_3\)-invariant.  

For the local side, in the case, 
\(\NN_L(P)/P\cong SD_{16}=\langle u,v\rangle\) by \eqref{p3:eq:SDmat}. For two fixed linear characters  
$\lambda_{+,+}$ and $\lambda_{+,-}$ of $D_8$, they extend to $SD_{16}$ by
\[
 \widehat\lambda_{+,+}(u)=1,\ \widehat\lambda_{+,+}(v)=1,
 \quad
 \widehat\lambda_{+,-}(u)=1,\ \widehat\lambda_{+,-}(v)=-1.
\]
These extensions are rational-valued. For the degree-$2$ character \(\rho\), we take \(\kappa\in\Irr(\langle u\rangle)\) with \(\kappa(u)=\zeta_8\), and let
 $\widehat\rho=\Ind_{\langle u\rangle}^{SD_{16}}(\kappa)$. 
The stabilizer of \(\kappa\) is \(\langle u\rangle\), so
\(\widehat\rho\) is irreducible of degree $2$. Its restriction to \(D_8=\langle u^2,v\rangle\) is \(\rho\). Since \(\kappa^{\sigma_3}=\kappa^3=\kappa^v\), induction of $\kappa$ 
gives \(\widehat\rho^{\sigma_3}=\widehat\rho\).
Thus, the three local extensions, inflated to \(\NN_L(P)\), are \(\HH_3\)-invariant.

For these pairs in the former three rows of Table \ref{tab:6}, \(\AT=A=L\). Conjugation by \(\NN_L(P)\) fixes the global and local extensions. Together with Galois invariance this makes them stable under the full combined stabilizer, as required.

Second, for every member of 
\( \{\varphi_2,\varphi_3,\varphi_7,\varphi_8,
\varphi_{16},\varphi_{17}\}\) in the latter three rows of Table \ref{tab:6}, it's \(\HH_3\)-invariant and has an
outer automorphism orbit of length two by Lemma \ref{equ}. Thus \(\AT=A=S\). The corresponding weight has the same stabilizer by Proposition~\ref{p3:prop:equivariant}.
Take \(\widehat\eta=\eta\) and \(\widehat\psi=\psi^\circ\), as required.

Finally, for the nine defect-zero characters, it is  obvious.
We complete the proof.
\end{proof}

\subsection{Verification at the prime \texorpdfstring{$p=5$}{p=5}}\label{sec:p5}
\subparagraph{(i) Global characters and two actions.}
All ordinary characters \(\chi_i\) and Brauer characters \(\varphi_j\) below are those of \cite{T}. 
The decomposition matrix of the principal block is reproduced in \cite[Block.1]{T}. Values of ordinary characters and the class fusion map \(S\hookrightarrow L\) are taken from the
named table records in \cite{OrdData} and Appendix \ref{app:maps}. 

There are one positive-defect \(5\)-block, the principal block \(B_0\), and six defect-zero blocks.  Moreover, the latter have
\begin{equation}\label{p5:dz}
 \dz_5(S)=\{\chi_7,\chi_8,\chi_{14},\chi_{15},\chi_{18},\chi_{19}\}
\end{equation}
with degrees $300,325,650,675,1300,1300$.
Their restrictions to $5$-regular elements are
$\varphi_9$, $\varphi_{10}$, $\varphi_{17}$, $\varphi_{18}$, $\varphi_{19}$, $\varphi_{20}$.
Here, the labels and degrees are directly from the top table of \cite{T}. Moreover, by \cite[Block.1]{T},  $\IBr(B_0)$ has fourteen irreducible Brauer characters
\begin{equation*}%\label{p5:eq:principal-labels}
\IBr_5(B_0)=\{\varphi_1,\varphi_2,\varphi_3,\varphi_4,\varphi_5,\varphi_6,
\varphi_7,\varphi_8,\varphi_{11},\varphi_{12},\varphi_{13},\varphi_{14},\varphi_{15},\varphi_{16}\}
\end{equation*}
with degrees $1,26,26,27,27,78,109,109,351,351,460,460,593,593$.
\begin{lemma}\label{lem:p=5 global characters}
All Brauer characters in defect-zero blocks are $(\HH_5\times \Out(S))$-invariant. Moreover, $\HH_5$ and $\Out(S)$ act respectively on $\IBr_5(B_0)$ by the non-trivial permutations 
\[
\tau
=(\varphi_2\ \varphi_3)(\varphi_7\ \varphi_8)
 (\varphi_{13}\ \varphi_{14})(\varphi_{15}\ \varphi_{16}),
\]
\[
\alpha=(\varphi_2\ \varphi_3)(\varphi_7\ \varphi_8)
 (\varphi_{13}\ \varphi_{14})(\varphi_{15}\ \varphi_{16})
 (\varphi_{19}\ \varphi_{20}),
\] 
and other characters are $(\HH_5\times\Out(S))$-invariant.
\end{lemma}
\begin{proof}
Let \(D\) be the  decomposition matrix of $B_0$ as in \cite[Block. 1]{T}. Then ${\rm rank}(D)=14$.
Restrict the sixteen ordinary rows in $B_0$ to the twenty \(5\)-regular classes, obtaining a matrix \(Y\). Write \(X\) for the table of Brauer characters for the principal block $B_0$. Then  $DX=Y$ and such 
$X$ is unique. The fifth-power map on classes computes $\sigma_5$ on the rows of $X$(see Appendix \ref{app:maps}).

On the other hand, we consider the class fusion map 
$S\hookrightarrow L$. Since $L/S\cong \Out(S)$,
two \(S\)-conjugacy classes \(c_1,c_2\) are fused in \(L\) if and only if $c_1=\alpha(c_2)$ for some $\alpha\in L/S$.
Thus this compute the action of the outer automorphism $\alpha\neq 1$. The input maps are written explicitly in  Appendix \ref{app:maps}. From \cite{T,OrdData}, the calculation gives  $\tau
=(\varphi_2\ \varphi_3)(\varphi_7\ \varphi_8)(\varphi_{13}\ \varphi_{14})(\varphi_{15}\ \varphi_{16})$ and $\alpha=(\varphi_2\ \varphi_3)(\varphi_7\ \varphi_8)(\varphi_{13}\ \varphi_{14})(\varphi_{15}\ \varphi_{16})(\varphi_{19}\ \varphi_{20})$, and other characters are fixed. 
Thus the two actions agree on $\IBr(B_0)$, with fixed set
$\{\varphi_1, \varphi_4,\varphi_5,\varphi_{11},\varphi_{12} \}$. They fix all six defect-zero characters under $\tau$, whereas $\alpha$ exchanges the two of degree $1300$.
We complete the proof.
\end{proof}


\subparagraph{(ii) Local characters and two actions.}
Let \(P:=P_5\in\Syl_5(S)\). Then $P\cong C_5\times C_5$, $\NN_S(P)\cong P\rtimes 4A_4$ and $\NN_L(P)\cong P\rtimes 4S_4$(see \cite{T}). 
\begin{lemma}\label{p5:prop:radical}
There are one non-trivial \(S\)-conjugacy classes of \(5\)-radical subgroups, represented by $P$.
\end{lemma}
\begin{proof}
Note that \(\OO_5(\NN_S(P))=P\) because \(P\) is a normal Sylow subgroup of its normalizer. 
Let \(Q=\langle z\rangle\leq P\) have
order \(5\). Since \(P\) is abelian, $P\leq\CC_S(Q)=\CC_S(z)$ which is of order $50$ because the ATLAS conjugacy-class table in \cite{A} has one class of elements of order $5$, labelled $5A$, with centralizer order $50$. Also, a group of order $50$ has a unique Sylow $5$-subgroup, so
\(\OO_5(\CC_S(Q))=P\). The centralizer is normal in \(\NN_S(Q)\), and then its \(5\)-core is characteristic. 
Thus $\OO_5(\NN_S(Q))=P>Q$. So \(Q\) is not radical. 
We are done.
\end{proof}

Therefore, every weight character comes either from $\dz_5(S)$ with trivial radical $5$-subgroup, or from
$\Irr(\NN_S(P)/P)$ with non-trivial radical $5$-subgroup $P$. Since $\NN_S(P)/P$ is a $5'$-group, all its irreducible characters are $5$-defect zero.


Let \(E=\NN_S(P)/P\) and \(\et=\NN_L(P)/P\). 
Then $E\cong 4.A_4$ and $\et\cong 4.S_4$.  In the standard affine realization, we obtain
\begin{equation}\label{p5:eq:central-product}
 E\cong C_4\circ\SL_2(3),\
 \ZZ(E)=\ZZ(\et)\cong C_4,\
 E/\ZZ(E)\cong A_4,\
 \et/\ZZ(\et)\cong S_4.
\end{equation}
The central product identifies the two central involutions.
The table-library entries in \cite{T} identify the affine normalizer $\NN_S(P)$  with \texttt{PrimitiveGroup(25,18)}, and $\NN_L(P)$  with \texttt{PrimitiveGroup(25,19)}, respectively.
Record the local inclusion through a  class fusion
in \cite{T}.


For describe the weight characters of $S$, we need some group structures as follows. In \(\GL_2(5)\), we take two elements
\[
a=\begin{pmatrix}2&0\\0&3\end{pmatrix},\quad b=\begin{pmatrix}0&1\\4&0\end{pmatrix},
\]
and write $K=\langle a,b\rangle\cong Q_8$.
Then the affine realizations above can be taken as
\begin{equation*}%\label{p5:eq:matrix-model}
 \et=\NN_{\GL_2(5)}(K),\quad
 E=\{g\in\et|\det(g)\in\{1,4\}\}.
\end{equation*}
This is the pair of local actions on \(P\cong\FF_5^2\) used below. For example, let
\[
t=\begin{pmatrix}3&3\\4&1\end{pmatrix},\quad
c=2I_2,\quad
u=\begin{pmatrix}0&4\\3&0\end{pmatrix}.
\]
Then \(t\) has order \(3\), \(c\) has order \(4\), and
$\langle a,b,t\rangle\cong\SL_2(3)$. Thus
$E=\langle a,b,t,c\rangle$ and  $\et=\langle E,u\rangle$.
Here \(c\) is central and \(u\notin E\) because \(\det(u)=3\). The projective images of \(E\) and \(\et\) are \(A_4\) and \(S_4\). These matrix calculations establish the descriptions in \eqref{p5:eq:central-product}.

Since \(5\nmid|\et|\), ordinary irreducible characters and Brauer characters of \(E\) and \(\et\) are naturally identified. 
Let $V$ be the natural two-dimensional \(\mathbb F_5E\)-module. Its second symmetric $\Sym^2 V$ is a three-dimensional \(\FF_5E\)-module. Write \(\nu\) for the Brauer character of $V$, regarded also as the corresponding ordinary character. 
Let
\[
\delta=\det|_E,\quad \theta=\operatorname{Br}(\Sym^2 V).
\]
Here, $\delta$ is viewed as a Brauer linear character by lifting its eigenvalue to a root of unity, and it has order \(2\).  $\theta$ denotes the Brauer character corresponding to the \(\FF_5E\)-module $\Sym^2 V$. Finally, choose a linear character \(\lambda\) of order \(3\) with the kernel $\langle K,c\rangle$.
\begin{lemma}\label{p5:lem:localchars}
The following list consists of all irreducible characters of \(E\)
\begin{equation*}%\label{p5:eq:localchars}
 \begin{split}
 \Irr(E)={}&\{\delta^e\lambda^j|\ e\in\{0,1\},\ j\in\{0,1,2\}\}\\
 &\sqcup\{\nu\delta^e\lambda^j| \ e\in\{0,1\},\ j\in\{0,1,2\}\}\\
 &\sqcup\{\theta,\theta\delta\},
 \end{split}
\end{equation*}
and their degrees are \(1^6,2^6,3^2\), where the superscripts denote multiplicities.
\end{lemma}
\begin{proof}
We claim that \(V\) is absolutely  irreducible on \(K\). Take the standard basis 
$e_1,e_2$ of $V$. Then $ae_1=2e_1$, $ae_2=3e_2$, $be_1=4e_2$ and  $be_3=e_1$.
Suppose there exists a one-dimensional \(K\)-invariant subspace \(W\). Since \(W\) is \(a\)-invariant, and the two eigenvalues \(2,3\) of \(a\) are distinct, \(W\) must be one of  $\langle e_1\rangle$ and $\langle e_2\rangle$. However,
$b\langle e_1\rangle=\langle e_2\rangle$ and 
$b\langle e_2\rangle=\langle e_1\rangle$. This shows that two one-dimensional subspace is not \(K\)-invariant. This is a contradiction. Thus 
$V$ is  irreducible on \(K\). Note that this argument still holds over the algebraic closure, so $V$ remains absolutely irreducible, as claimed.
Since \(K\leq E\), $\nu\in \Irr(E)$ with $\nu(1)=2$,
and thus $\nu\delta^e\lambda^j\in\Irr(E)$.

Note that $\dim \Sym^2V=3$, and then take the basis $e_1^2, e_1e_2, e_2^2$. By the action of $a$ and $b$, we have three common eigenvectors
$v_1=e_1^2+e_2^2$, $v_2=e_1^2-e_2^2$, $v_3=e_1e_2$.
Thus ${\rm Res}_K^E(\Sym^2V)=\bigoplus_{i=1}^3 \langle v_i \rangle$. Since $K/\ZZ(K)\cong C_2\times C_2$, the restriction ${\rm Res}_K^E(\theta)$ is the sum of three non-trivial linear characters $\alpha_i$.  Since $tat^{-1}=b^{-1}$ and $tbt^{-1}=-ab$, \(t\)-conjugation cycles these three $\alpha_i$,
Suppose that  $0\neq W\subseteq \Sym^2V$ is an 
$E$-invariant subspace. Since $K\unlhd E$ and $5\nmid|K|$, $W$ is reducible on $K$.
Note that $W$ must be the direct sum of several of them. Due to the cyclic permutation of these three constituents, there can only be $W=\Sym^2V$.
Thus $\theta\in \Irr(E)$ with $\theta(1)=3$.
The same is true for $\theta\delta$.

Since $c=2I_2$, \(c\) is a scalar matrix on \(V \).
Under the lift of Brauer characters,  
$(\nu\delta^e\lambda^j)(c)=(-1)^e\nu(c)$ because $\delta(c)=-1$, $\lambda(c)=1$ and $\nu(c)\neq0$. On the other hand, the trace of $t$ is $-1$
and  its determinant is $1$, so the characteristic polynomial of $t$ is $x^2+x+1$. Its two eigenvalues, when lifted to the complex numbers, are $\omega,\omega^2$. Thus $\nu(t)=\omega+\omega^2=-1$.
Also, since $\delta(t)=1$ and $\lambda(t)=\omega$,
we obtain $\nu\delta^e\lambda^j)(t) =-\omega^j$.
Therefore, the six degree-$2$ characters $\nu\delta^e\lambda^j$ are pairwise distinct.
The six linear characters $\delta^e\lambda^j$ are distinct since \(\delta\) and \(\lambda\) have coprime orders. 

Finally, 
note that $6\cdot1^2+6\cdot2^2+2\cdot3^2=48=|E|$, so there cannot be any other irreducible characters.
\end{proof}

We consider now the actions of outer auromorphisms and Galois automorphisms on the set $\Irr(E)$.
\begin{lemma}\label{lem:p=5 local characters}
The Brauer restrictions of these six characters $1,\delta,\nu,\nu\delta,\theta,\theta\delta$ are $(\HH_5\times(\et\setminus E))$-invariant. 
Moreover, the action of $\HH_5\times(\et\setminus E)$ on other characters induced by the map $\lambda\mapsto\lambda^{-1}$.
\end{lemma}
\begin{proof}
Note that the six  characters $1,\delta,\nu,\nu\delta,\theta,\theta\delta$ in Lemma \ref{p5:lem:localchars} are afforded by \(\FF_5\)-modules  $1,\delta, V, V\otimes \delta, \Sym^2 V, \Sym^2 V\otimes \delta$, respectively. Then their Brauer values are fixed by 
\(\sigma_5\in\HH_5\).

Note that conjugation by an element of \(\et\setminus E\) fixes the six characters. On the other hand, $\lambda^{\sigma_5}=\lambda^5=\lambda^{-1}$ because 
$\lambda$ has order $3$. An odd permutation in \(S_4\) inverts \(A_4/V_4\cong C_3\), so $\lambda^\alpha=\lambda^{-1}$ for any $\alpha\in \et\setminus E$. Also, $\delta$ and $\nu$ have order $2$. Therefore, we obtain that
 $(\delta^e\lambda^j)^{\sigma_5}=(\delta^e\lambda^j)^\alpha=\delta^e\lambda^{-j}$ and
 $(\nu\delta^e\lambda^j)^{\sigma_5}=(\nu\delta^e\lambda^j)^\alpha=\nu\delta^e\lambda^{-j}$.
We complete the proof.
\end{proof}

Thus, for all members of $\Irr(E)$, both actions have
six  $\sigma_5\alpha$-fixed points and four 
$\sigma_5\alpha$-orbits. Together with Lemma \ref{lem:p=5 global characters}, we have the following correspondence.
\begin{table}[htbp]
\centering
\caption{The correspondence between Brauer characters and weight characters}\label{tab:7}
\begin{tabular}{lll}
\toprule
Brauer characters & Degrees & Weight characters\\
\midrule
\(\varphi_1\) & 1 & \(1\)\\
\(\varphi_6\) & 78 & \(\delta\)\\
\(\varphi_4\) & 27 & \(\nu\)\\
\(\varphi_5\) & 27 & \(\nu\delta\)\\
\(\varphi_{11}\) & 351 & \(\theta\)\\
\(\varphi_{12}\) & 351 & \(\theta\delta\)\\
\((\varphi_2,\varphi_3)\) & \((26,26)\) & \((\lambda,\lambda^2)\)\\
\((\varphi_7,\varphi_8)\) & \((109,109)\) & \((\delta\lambda,\delta\lambda^2)\)\\
\((\varphi_{13},\varphi_{14})\) & \((460,460)\) & \((\nu\lambda,\nu\lambda^2)\)\\
\((\varphi_{15},\varphi_{16})\) & \((593,593)\) & \((\nu\delta\lambda,\nu\delta\lambda^2)\)\\
\bottomrule
\end{tabular}
\end{table}

\begin{proposition}\label{p5:prop:equivariant}
There exists an 
\((\HH_5\times\Aut(S))\)-equivalent bijection
\(\IBr_5(S)\longrightarrow\mathcal W_5(S)\).
\end{proposition}
\begin{proof}
Define \(\Omega\) on the defect-zero blocks by
$\chi^\circ\mapsto\overline{(1,\chi)}$ for every $\chi\in\dz_5(S)$ in \eqref{p5:dz}. For the principal block $B_0$, we choose the correspondence in Table \ref{tab:7}. We compare both actions in Lemma \ref{lem:p=5 global characters} with in Lemma \ref{lem:p=5 local characters},  both sides fix the first--six rows and interchange the entries in each of the four paired rows. The inner automorphisms act trivially on both sets of conjugacy classes. We are done.
\end{proof}


\subparagraph{(iii) Compatible extensions.} The following verifications concern  extensions of Brauer characters and weight characters for the above bijection. Choose
\(\AT=L_{\eta^{\HH_5}}\) and \(A=L_\eta\) as in Proposition \ref{prop:criterion}.
\begin{proposition}\label{thm:p5}
Let $\eta\in\IBr_5(S)$ and $\Omega(\eta)=\overline{(R,\psi)}\in\WW_5(S)$.
Then $\eta$ and $\psi^\circ$ satisfy the condition \eqref{extension criterion}.
\end{proposition}
\begin{proof}
For the Sylow $5$-subgroup \(P\), the Frattini argument gives \(L=S\NN_L(P)\), so $L/S\cong \NN_L(P)/\NN_S(P)$, as required. 
Note that $\eta$ comes from restriction of $\dz_5(S)$ and $\IBr_5(B_0)$, and that $\psi$ comes from Lemma \ref{p5:lem:localchars}. 

First, for extensions of the fixed characters
$\varphi_1,\varphi_4,\varphi_5,\varphi_6,
\varphi_{11},\varphi_{12}\in\IBr_5(B_0)$. 
As in \cite[Block.1]{T}, 
$\varphi_j=\chi_j^\circ\ (j=1,4,5,6)$, 
$\varphi_{11}=\chi_{10}^\circ$ and $\varphi_{12}=\chi_{11}^\circ$. For these six ordinary characters $\chi_j\ (j=1,4,5,6)$, $\chi_{10}$, $\chi_{11}$, there are extensions to \(L\) with values in \(\QQ(i)\), 
where $i$ is the imaginary unit.  More precisely, 
the following rows of $L$  in the record \texttt{MOT("2F4(2)'.2",...)} of \cite{OrdData} are suitable as follows.
\begin{center}
\begin{tabular}{cccc}
\toprule
Brauer characters of \(S\) & Ordinary characters of \(S\)& Row in \(L\)& Degrees\\
\midrule
\(\varphi_1\)&\(\chi_1\)&1&1\\
\(\varphi_4\)&\(\chi_4\)&4&27\\
\(\varphi_5\)&\(\chi_5\)&6&27\\
\(\varphi_6\)&\(\chi_6\)&8&78\\
\(\varphi_{11}\)&\(\chi_{10}\)&16&351\\
\(\varphi_{12}\)&\(\chi_{11}\)&18&351\\
\bottomrule
\end{tabular}
\end{center}
Here, these row numbers in $L$ refer to the ordinary character table of \(L\) record. By their restrictions to $S$, we obtain characters of \(S\). 


Let \(\widehat\chi\) be any of extensions of the six chosen ordinary characters, and \(\eta={\rm Res}_S^L(\widehat\chi^\circ)\). 
By an ordinary representation of \(L\) and reduction modulo \(5\), this gives a $\FF_5$-representation of \(L\) whose restriction to \(S\) has the character \(\eta\), which is irreducible. Note that the above  Brauer restriction
is a simple module, so the representation of \(L\) is itself simple. Thus \(\widehat\chi^\circ\) is an extension of \(\eta\).
By the decomposition matrix of \(L\)(see \cite[Block~1, pp.1--2]{T}),  independently confirm these six characters.
Since \(i^5=i\), every such extension is \(\HH_5\)-invariant. In particular, all $\varphi_1,\varphi_4,\varphi_5,\varphi_6,\varphi_{11},\varphi_{12}$ have an \(\HH_5\)-invariant extension to \(L\), denoted by 
\begin{equation}\label{p5:eq:global extensions}
\widehat{\varphi_1}, \widehat{\varphi_4}, \widehat{\varphi_5},
\widehat{\varphi_6}, \widehat{\varphi_{11}}, \widehat{\varphi_{12}}    
\end{equation}

For the fixed characters $1,\delta,\nu,\nu\delta,\theta,\theta\delta\in\Irr(E)$, we will construct their extensions. It is easy to check that each also extends to \(\et\) by the corresponding module extensions. Using its natural module \(\widehat V=\FF_5^2\), its character \(\widehat\delta=\det|_{\widetilde E}\), and the six extended modules 
\begin{equation}\label{p5:eq:local extensions}
 1,\ \widehat\delta,\ \widehat V,\ \widehat V\otimes\widehat\delta,\
 \Sym^2\widehat V,\ \Sym^2\widehat V\otimes\widehat\delta.
\end{equation}
Since the six extended modules are defined over \(\FF_5\), they are also \(\HH_5\)-invariant. We choose \(\AT=A=L\). Match global extensions in \eqref{p5:eq:global extensions} with  local extensions in \eqref{p5:eq:local extensions}, inflated from \(\et\) to \(\NN_L(P)\). Therefore, both extensions are \((\HH_5\times \NN_L(P))\)-invariant by Lemmas \ref{lem:p=5 global characters} and \ref{lem:p=5 local characters}.

Second, let \(\eta\) be one of the last four rows in Table \ref{tab:7}. Then by Lemma \ref{lem:p=5 global characters}, $\eta^\alpha=\eta^{\sigma_5}\ne\eta$. It is easy to check that $L_\eta=S$ and $L_{\eta^{\HH_5}}=L$. Thus \(\AT=L\), but \(A=S\). In particular, an element combining an outer automorphism with an odd Frobenius power can fix $\eta$.
On the other hand, note that $A/S=1=\NN_A(P)/\NN_S(P)$, so
there is no non-trivial extension, as required.

Finally, for the six defect-zero characters, $L_{\chi_i}=L$ for $i=7,8,14,15$, and $L_{\chi_j}=S$ for $j=18,19$. Choosing identical extensions on the two sides,
the conclusion holds. We complete the proof.
\end{proof}
\begin{theorem}\label{thm: all p}
The simple groups ${}^2G_2(3)'$ and ${}^2F_4(2)'$ satisfy the inductive NAW condition at all primes $p$.
\end{theorem}
\begin{proof}
Since ${}^2G_2(3)'\cong\operatorname{PSL}_2(8)$, this case is handled by \cite[Proposition 10.8]{FFZ}. Thus we suppose that $S={}^2F_4(2)'$. In the case, we have $p=2,3,5,13$. For $p=2$, $S$ satisfies the condition by \cite[Proposition 9.3]{FFZ}. If $p=13$, then by the order $|S|$, every Sylow $13$-subgroup of $S$ is cyclic. Thus we apply Proposition \ref{cyclic Sylow}. 
For $p=3$, Propositions \ref{p3:prop:equivariant} and \ref{thm:p3} show that Proposition \ref{prop:criterion} holds.  Propositions \ref{p5:prop:equivariant} and \ref{thm:p5} verify the case of $p=5$. We are done.
\end{proof}

\section{Proof of main theorem}\label{section 5}
\begin{proof}[Proof of Theorem \ref{thm:main}]
For the simple groups ${}^2B_2(2^{2m+1})$ or ${}^2G_2(3^{2m+1})$ with $m\geq1$, it follows from
Theorem \ref{thm:Suzuki} that they satisfy the inductive NAW condition at all primes $p$. 

For the simple ${}^2F_4(2^{2m+1})$ with $m\geq 1$, \cite[Theorem C]{FFZ} has already addressed  the case of $p=2$. Theorem \ref{thm:2F4-p=3} and Theorem \ref{thm:2F_4 for p>=5} respectively verify the cases of $p=3$ and $p\geq 5$.  

For the simple groups ${}^2F_4(2)'$ and ${}^2G_2(3)'$, Theorem \ref{thm: all p} proves that they satisfy the inductive NAW condition at all primes $p$. We complete the proof.
\end{proof}

\appendix
\section{Appendix: Class indexing, power maps and fusion maps}\label{app:maps}
When \(S={}^2F_4(2)'\) and \(L={}^2F_4(2)\). The twenty-two ordinary  class positions of \(S\) in \cite{OrdData} are
\[
\begin{gathered}
(1A,2A,2B,3A,4A,4B,4C,5A,6A,8A,8B,8C,8D,
10A,12A,12B,13A,13B,16A,16B,16C,16D).
\end{gathered}
\]
For \(p=3\), delete positions \(4,9,15,16\); for \(p=5\), delete positions \(8,14\). Thus the $p$-regular class counts are \(18\) and \(20\), respectively.

The third-power and fifth-power maps, given by their images on the full list of positions, are
\begin{align*}
 \pi_3={}&(1,2,3,1,5,6,7,8,3,11,10,12,13,14,5,5,17,18,22,21,20,19),\\
 \pi_5={}&(1,2,3,4,5,6,7,1,9,10,11,12,13,2,16,15,18,17,21,22,19,20).
\end{align*}
These are the third and fifth power-map fields of the ordinary \(S\)-record. The class fusion map of $S$ into the twenty-nine classes of $L$ is
\[
(1,2,3,4,5,6,7,8,9,10,11,12,12,13,14,14,15,15,16,17,16,17).
\]
It is the \texttt{ALF} record from \(S\) to \(L\) in \cite{OrdData} and gives the permutation of outer automorphisms
\[
\alpha=(12\ 13)(15\ 16)(17\ 18)(19\ 21)(20\ 22).
\]
For any Brauer row \(b\), \(b^{\sigma_p}(c)=b(c^p)\). Apply this identity on the $p$-regular classes to \(X_p\), after solving \(D_pX_p=Y_p\), and apply the  outer automorphism permutation of conjugate class in the same way. 



\phantomsection
\addcontentsline{toc}{section}{References}
\begin{thebibliography}{99}
\bibitem{A} 
R. Abbott, J. Bray, S. Linton, S. Nickerson, S. Norton, R. Parker, I. Suleiman, J. Tripp, P. Walsh, and R. Wilson, \emph{Atlas of finite group representations--version 3}. \url{https://brauer.maths.qmul.ac.uk}.

\bibitem{Alperin1987}
J. L. Alperin, \emph{Weights for finite groups}. In: The Arcata Conference on Representations of Finite Groups, Arcata, Calif. (1986), Part I. Proc. Sympos. Pure Math., vol. {\bf 47}, pp. 369--379, Am. Math. Soc., Providence, 1987.

\bibitem{An}
J. An, \emph{The Alperin and Dade conjectures for Ree groups ${}^2F_4(q^2)$ in non-defining characteristics}, J. Algebra \textbf{203}(1) (1998), 30--49.
%\url{https://doi.org/10.1006/jabr.1997.7287}.


\bibitem{OrdData}
T. Breuer, \emph{GAP ordinary character-table data file \texttt{ctotwis1.tbl}}, header revision 1.1. \url{https://raw.githubusercontent.com/jmichel7/gap3-jm/master/tbl/ctotwis1.tbl}. 

\bibitem{T}
T. Breuer, \emph{The GAP character table library, Version 1.3.11}, 2025, \url{https://www.math.rwth-aachen.de/~Thomas.Breuer/ctbllib/}.


\bibitem{DL}
P. Deligne and G. Lusztig, \emph{Representations of reductive groups over finite fields}, Ann. of Math. {\bf 103}(1) (1976), 103--161. 
%\url{https://doi.org/10.2307/1971021}.

\bibitem{DM}
F. Digne and J. Michel, \emph{Complements on disconnected reductive groups},  Pacific J. of Math. \textbf{279} (2015), 203--228. 
%\url{doi: 10.2140/pjm.2015.279.203}.

\bibitem{FFZ}
Z. Feng, Q. Fu and Y. Zhou, \emph{A reduction theorem for the Navarro Alperin weight conjecture}, arXiv:2312.02594v5, 2026. 
\url{https://arxiv.org/abs/2312.02594v5}.

\bibitem{Ghaffarzadeh}
M. Ghaffarzadeh, \emph{Character degrees of extensions of the Suzuki groups \({}^2B_2(q^2)\)}, J. Algebra Appl. {\bf 17}(1) 1850006 (2018). 
%\url{https://doi.org/10.1142/S0219498818500068}

\bibitem{GM20}
M. Geck and G. Malle, \emph{The Character Theory of Finite Groups of Lie Type: A Guided Tour}, Cambridge Studies in Advanced Mathematics, vol. {\bf 187}, Cambridge University Press, Cambridge, 2020.


\bibitem{GH}
G. Hiss and M. Geck, \emph{Basic sets of Brauer characters of finite groups of Lie type}, J. Reine Angew. Math.
{\bf 418} (1991), 173--188. %\url{https://doi.org/10.1515/crll.1991.418.173}.


\bibitem{Himstedt}
F. Himstedt, \emph{On the decomposition numbers of the Ree groups ${}^2F_4(q^2)$ in non-defining character
istic}, J. Algebra \textbf{325}(1) (2011), 364--403.
%\url{https://doi.org/10.1016/j.jalgebra.2010.08.030}.

\bibitem{HH}
F. Himstedt and S. C. Huang, \emph{Dade's invariant conjecture for the Ree groups ${}^2F_4(q^2)$ in defining characteristic}, Comm. Algebra, \textbf{40}(2) (2012), 452--496.
%\url{https://doi.org/10.1080/00927872.2010.531994}.

\bibitem{HHM}
G. Hiss, W. J. Husen, and K. Magaard, \emph{Imprimitive irreducible modules for finite quasisimple groups}, Mem. Amer. Math. Soc. \textbf{234}(1104) (2015).
%\url{https://doi.org/10.1090/memo/1104}.

\bibitem{I2007}
I.M. Isaacs, G. Malle and G. Navarro, \emph{A reduction theorem for the McKay conjecture}, Invent. Math. \textbf{170} (2007), 33--101.

\bibitem{Johansson}
B. Johansson, \emph{The inductive McKay-Navarro condition for the Suzuki and Ree groups and for groups with non-generic Sylow normalizers}, J. Pure Appl. Algebra \textbf{228}(2)  (2024), 107469. %\url{https://doi.org/10.1016/j.jpaa.2023.107469}.


\bibitem{Malle}
G. Malle, \emph{On the inductive Alperin-McKay and Alperin weight conjecture for groups with abelian Sylow subgroups},
J. Algebra \textbf{397} (2014), 190--208.
%\url{https://doi.org/10.1016/j.jalgebra.2013.09.013}.

\bibitem{MT18}
K. Muraleedaran and D.E. Taylor, \emph{Normalisers of parabolic subgroups in finite unitary reflection groups}, 
J. Algebra \textbf{504} (2018), 479--505.
%\url{https://doi.org/10.1016/j.jalgebra.2018.02.026}.


\bibitem{Navarro2004}
G. Navarro, \emph{The Mckay conjecture and Galois automorphisms}, Ann. of Math. {\bf 160}(2) (2004), 1129--1140.
%\url{https://doi.org/10.4007/annals.2004.160.1129}.

\bibitem{Orevkov}
 S. Yu. Orevkov. \emph{Products of conjugacy classes in finite unitary groups ${\rm GU}(3,q^2)$ and ${\rm SU}(3,q^2)$},
 Ann. Fac. Sci. Toulouse Math. \textbf{22}(2) (2013), 219--251.
%\url{doi:10.5802/afst.1371}.

\bibitem{T2013}
A. Turull, \emph{The strengthened Alperin weight conjecture for $p$-solvable groups}, J. Algebra \textbf{398} (2014), 469--480.

\bibitem{S16}
I. Schwabrow. \emph{The centre of a block}. PhD Thesis, University of Manchester, 2016.


\bibitem{W13}
D. L. White, \emph{Character degrees of extensions of
\(\mathrm{PSL}_2(q)\) and \(\mathrm{SL}_2(q)\)},
J. Group Theory. \textbf{16} (2013), 1--33.
%\url{https://doi.org/10.1515/jgt-2012-0026}.



\bibitem{R}
R. Wilson, J. Thackray, R. Parker, F. Noeske, J. Müller, F. Lübeck, C. Jansen, G. Hiss, and T. Breuer, \emph{The Modular Atlas Project}, \url{https://www.math.rwth-aachen.de/homes/MOC/decomposition/}.



%\bibitem{FFZGeneral} Zhicheng Feng, Qulei Fu, and Yuanyang Zhou. \emph{A reduction theorem for the blockwise Navarro Alperin weight conjecture via H-triples}, 2026. arxiv.org/abs/2512.15243v2.

%\bibitem{NavarroTiep} G.~Navarro and P.~H.~Tiep, \emph{A reduction theorem for the Alperin weight conjecture}, Invent. Math. \textbf{184} (2011), 529--565. \url{https://doi.org/10.1007/s00222-010-0295-2}. 

\end{thebibliography}
\end{document}
